
\begin{enumerate}
\def\labelenumi{\arabic{enumi}.}
\setcounter{enumi}{17}
\tightlist
\item
  Let \(G\) be a connected real Lie group, \(Z\) its centre, \(Z_0\) the identity component of \(Z\) , \(g\) the Lie algebra of \(G\) and \(h\) the centre of \(g\) . Then \(Z\) and \(Z_0\) are Lie quasi-subgroups of \(G\) with Lie algebra \(h\) (Exercise 3). An acceptable norm on \(g\) is a norm defining the topology of \(g\) and making \(g\) into a normed Lie algebra.
\end{enumerate}

\begin{enumerate}
\def\labelenumi{(\alph{enumi})}
\item
  For all \(r > 0\) and \(z \in \mathfrak{z}\) , there exists an acceptable norm on \(\mathfrak{g}\) whose value at \(z\) is \(< r\) . (Let \(q\) be an acceptable norm on \(\mathfrak{g}\) . Show that, for suitably chosen \(\lambda > 0\) , the function \(x \mapsto \|x\| = \lambda q(x) + \inf_{y \in \mathfrak{z}} q(x + y)\) has the required properties.)
\item
  Show that the following conditions are equivalent:
\item
  \(Z_{0}\) is simply connected;
\end{enumerate}

\begin{enumerate}
\def\labelenumi{(\roman{enumi})}
\setcounter{enumi}{1}
\item
  for every acceptable norm on g, the restriction of \(\exp_{G}\) to the open ball of centre 0 and radius \(\pi\) is injective;
\item
  there exists \(r > 0\) such that, for every acceptable norm on \(\mathfrak{g}\) , the restriction of \(\exp_{\mathbb{G}}\) to the open ball of centre 0 and radius \(r\) is injective.
\end{enumerate}

(To prove (iii) \(\Rightarrow\) (i), use (a). To prove (i) \(\Rightarrow\) (ii), suppose that \(x, y\) are in \(\mathfrak{g}\) , \(\| x \| < \pi\) , \(\| y \| < \pi\) , \(x \neq y\) , \(\exp_{\mathfrak{g}} x = \exp_{\mathfrak{g}} y\) . Then \(\exp \mathrm{d} x = \exp \mathrm{d} y\) and hence \(\mathrm{ad} x = \mathrm{ad} y\) by Proposition 17 of §6 applied to a complexification of \(\mathfrak{g}\) . Hence there exists a non-zero \(z\) in \(\mathfrak{z}\) such that \(\exp_{\mathfrak{g}} z = e\) ; it follows that (i) is false.)

\begin{enumerate}
\def\labelenumi{(\alph{enumi})}
\setcounter{enumi}{2}
\tightlist
\item
  By considering the group \(\mathbf{G}'\) of Exercise 13 (b), show that the conclusion of (b) is no longer true if \(\pi\) is replaced by a number \(> \pi\) . (In the notation of Exercise 13, use the norm \(ae_{1} + be_{2} + ce_{3} \mapsto (a^{2} + b^{2} + c^{2})^{\frac{1}{2}}\) on \(\mathrm{L}(\mathbf{G}')\) .)
\end{enumerate}

\begin{enumerate}
\def\labelenumi{\arabic{enumi}.}
\setcounter{enumi}{18}
\item
  Let G be a locally compact group and H a closed subgroup of G. Suppose that H is a finite-dimensional simply connected solvable Lie group and that G/H is compact. Show that there exists a compact subgroup L of G such that G is the semi-direct product of L and H. (Argue by induction on dim H. Consider the last non-trivial derived group of H and use Integration, Chapter VII, § 3, Proposition 3.)
\item
  Let \(G\) be a finite-dimensional solvable connected real Lie group. We introduce the notation \(S, S', F, \sigma\) of §6, no. 10, proof of Proposition 20.
\end{enumerate}

\begin{enumerate}
\def\labelenumi{(\alph{enumi})}
\item
  Using Proposition 21, show that \(\sigma\) is an isomorphism of S onto a Lie subgroup of the underlying real Lie group of S'.
\item
  Deduce that the universal complexification \(\tilde{G}\) of \(G\) is identified with \(S' / \sigma(F)\) and that the canonical mapping of \(G\) into \(\tilde{G}\) is an isomorphism of \(G\) onto a Lie subgroup of the underlying real Lie group of \(\tilde{G}\) .
\end{enumerate}

¶ 21. (a) Let G be a simply connected solvable real Lie group with the following properties: (α) L(G), which is n-dimensional, has a commutative ideal of dimension n - 1, corresponding to a subgroup A of G; (β) there exists an element σ of the centre of G which does not belong to A. Show that there exists an element \(x\) of \(\mathrm{L}(\mathbf{G})\) such that \(\exp x = \sigma\) . Show that \(\mathrm{L}(\mathbf{G})\) is the product of a commutative ideal and an ideal with a basis

\[
(x, a _ {1}, b _ {1}, \dots , a _ {k}, b _ {k})
\]

such that \(a_1, b_1, \ldots, a_k, b_k\) belong to \(\mathrm{L}(\mathrm{A})\) , \([x, a_i] = 2\pi n_i b_i\) , \([x, b_i] = -2\pi n_i a_i\) ( \(n_i \in \mathbf{Z} - \{0\}\) ) for all \(i\) . Generalize the results of Exercise 13 (d) to G.

\begin{enumerate}
\def\labelenumi{(\alph{enumi})}
\setcounter{enumi}{1}
\tightlist
\item
  Let G be a finite-dimensional simply connected solvable real Lie group. Let D be a discrete subgroup of the centre of G. There exist a basis \((x_{1}, x_{2}, \ldots, x_{n})\) of L(G) and an integer \(r \leqslant n\) with the following properties: (α) every element of G can be written uniquely in the form
\end{enumerate}

\[
\left(\exp t _ {1} x _ {1}\right) \dots \left(\exp t _ {n} x _ {n}\right)
\]

where \(t_1, \ldots, t_n\) are in \(\mathbf{R}\) ; \((\beta)x_1, \ldots, x_r\) are pairwise permutable and

\[
(\exp x _ {1}, \dots , \exp x _ {r})
\]

is a basis of the commutative group D. (Argue by induction on the dimension of G. Let a be maximal commutative ideal of L(G). Let A be the corresponding integral subgroup. Then A is closed in G and DA/A is a discrete subgroup of the centre of G/A, to which the induction hypothesis can be applied, whence there are elements \(x_1^*, \ldots, x_m^*\) of L(G/A) and an integer \(s\) . For \(1 \leqslant i \leqslant s\) , let \(\sigma_i\) be an element of D whose class modulo A is exp \(x_1^*\) . Construct one by one, using (a), representatives \(x_1, \ldots, x_m\) of \(x_1^*, \ldots, x_m^*\) such that exp \(x_i = \sigma_i\) and \([x_1, x_j] = 0\) for \(1 \leqslant i, j \leqslant s\) .

\begin{enumerate}
\def\labelenumi{(\alph{enumi})}
\setcounter{enumi}{2}
\tightlist
\item
  Deduce from (b) that every finite-dimensional connected solvable real Lie group is homomorphic to a space \(\mathbf{R}^n\times \mathbf{T}^m\) ( \(m,n\) integers \(\geqslant 0\) ).
\end{enumerate}

\begin{enumerate}
\def\labelenumi{\arabic{enumi}.}
\setcounter{enumi}{21}
\item
  Let \(G\) be a finite-dimensional connected real Lie group such that \(L(G)\) is reductive. Then \(G = (V \times S)/N\) where \(V\) is a finite-dimensional real vector space, where \(S\) is a simply connected semi-simple real Lie group and \(N\) is a central discrete subgroup of \(V \times S\) . Then \(G/\overline{D^1}G\) is isomorphic to \(V/\mathrm{pr}_1\overline{N}\) . Hence \(G/\overline{D^1}N\) is compact if and only if \(\mathrm{pr}_1\mathrm{N}\) generates the vector space \(V\) .
\item
  \begin{enumerate}
  \def\labelenumii{(\alph{enumii})}
  \tightlist
  \item
    Let G be a finite-dimensional commutative real Lie group with only a finite number of connected components. For G to be compact, it is necessary and sufficient that every finite-dimensional analytic linear representation of G on a complex vector space be semi-simple. (Use the proof of Proposition 32.)
  \end{enumerate}
\end{enumerate}

\begin{enumerate}
\def\labelenumi{(\alph{enumi})}
\setcounter{enumi}{1}
\tightlist
\item
  Let G be a finite-dimensional commutative complex Lie group with only a finite number of connected components. Let N be the kernel of \(\exp_{G}\) .
\end{enumerate}

For N to generate over C the vector space \(L(G)\) , it is necessary and sufficient that every finite-dimensional analytic linear representation of G be semi-simple. (Use (a), Lemma 1 and Proposition 32.)

\begin{enumerate}
\def\labelenumi{\arabic{enumi}.}
\setcounter{enumi}{23}
\item
  The Lie group \(\mathbf{SL}(2,\mathbf{R})\) is connected and almost simple, but \(\{I,-I\}\) is a normal commutative subgroup of \(\mathbf{SL}(2,\mathbf{R})\) .
\item
  Let G be an almost simple connected real or complex Lie group. Let A be a normal subgroup of G. If \(A \neq G\) , A is discrete and central. (Use Exercise 8 of §4.) Therefore the quotient of G by its centre is simple as an abstract group.
\item
  Let G be a finite-dimensional connected real Lie group. Suppose that G admits a finite-dimensional injective continuous linear representation \(\rho\) . Then (G, G) is closed in G. (Let R be the radical of G and S a maximal semi-simple integral subgroup of G. Using Exercise 7 (b), reduce it to the case where G is the semi-direct product of S and R. By Chapter I, § 6, Proposition 6, \(\rho\) is unipotent on (G, R). Hence \(\rho((G, R))\) is closed in the linear group and therefore (G, R) is closed in G.)
\item
  Let G be a finite-dimensional solvable simply connected real Lie group and N the largest connected nilpotent normal subgroup of G. Then G admits a finite-dimensional injective continuous linear representation which is unipotent on N. (Argue by induction on the dimension of G. Use Proposition 20 and Chapter I, § 7, Theorem 1.)†
\end{enumerate}

¶ 28. (a) Let k be a commutative field of characteristic 0, L an n-dimensional Lie algebra over k, \(L_{1}\) and \((n-1)\) -dimensional subalgebra containing no non-zero ideal of L and \(a_{0}\) an element of L not belonging to \(L_{1}\) . For \(i=2,3,\ldots\) , we define \(L_{i}\) inductively to be the set of \(x\in L_{i-1}\) such that \([x,a_{0}]\in L_{i-1}\) . We write \(L_{i}=L\) for \(i\leqslant0\) . Show that \([L_{0},L_{i}]\subset L_{i+j}\) (by induction on \(i+j\) ) and then that, for \(0\leqslant i\leqslant n\) , \(L_{i}\) is a subalgebra of L of codimension i (by induction on i).

\begin{enumerate}
\def\labelenumi{(\alph{enumi})}
\setcounter{enumi}{1}
\tightlist
\item
  For 0 \textless{} i \textless{} n, choose in \(L_{i}\) an element \(a_{i}\) not belonging to \(L_{i+1}\) such that \([a_{0}, a_{i}] \equiv ia_{i-1} \pmod{L_{i}}\) . Show, by induction on the ordered pairs \((i, j)\) ordered lexicographically, that, for \(0 \leqslant i < j < n\) , \(i + j - 1 < n\) and
\end{enumerate}

\[
[ a _ {i}, a _ {j} ] \equiv (j - i) a _ {i + j - 1} (\mathrm{mod}. \mathrm{L} _ {i + j}).
\]

\begin{enumerate}
\def\labelenumi{(\alph{enumi})}
\setcounter{enumi}{2}
\item
  Deduce that L is either of dimension 1, or of dimension 2 and noncommutative, or isomorphic to \(\mathfrak{sl}(2,k)\) .
\item
  Let A be a finite-dimensional real Lie algebra. A subalgebra B of A is called extendable if there exists a subalgebra \(B_{1} \supset B\) such that dim \(B_{1} = \dim B + 1\) . Let \(R'(A)\) denote the intersection of all the nonextendable subalgebras of A. Let \(R(A)\) denote the largest ideal of A with a composition series as an A-module all of whose quotients are of dimension 1.
\end{enumerate}

Show that \(\mathbf{R}'(\mathbf{A})\) is a characteristic ideal of \(\mathbf{A}\) . (Observe that \(\mathbf{R}'(\mathbf{A})\) is stable under \(\operatorname{Aut}(\mathbf{A})\) .)

\begin{enumerate}
\def\labelenumi{(\alph{enumi})}
\setcounter{enumi}{4}
\item
  Show that \(\mathbf{R}'(\mathbf{A})\supset \mathbf{R}(\mathbf{A})\) and that \(\mathrm{R}'(\mathrm{A} / \mathrm{R}(\mathrm{A})) = \mathrm{R}'(\mathrm{A}) / \mathrm{R}(\mathrm{A})\)
\item
  Show that, if B is a subalgebra of A, then R'(B) ⊃ R'(A) ∩ B.
\item
  If A is solvable, \(\mathrm{R}^{\prime}(\mathrm{A}) = \mathrm{R}(\mathrm{A})\) . (It can be assumed that \(\mathrm{R}(\mathrm{A}) = \{0\}\) , by (e). Suppose that \(\mathrm{R}^{\prime}(\mathrm{A}) \neq \{0\}\) . By (d), there exists a minimal non-zero ideal I of A contained in \(\mathrm{R}^{\prime}(\mathrm{A})\) . Using (f) and Chapter I, §5, Corollary 1 to Theorem 1, show that dim I = 1, whence a contradiction.)
\item
  Let P be the radical of \(\mathrm{R}^{\prime}(\mathrm{A})\) , B a semi-simple subalgebra of A and \(C = P + B\) which is a subalgebra of A by (d). Using (f), show that \(ad_{P}B\) can be expressed in a triangular form with respect to a suitable basis of P. Conclude that \([P, B] = \{0\}\) .
\end{enumerate}

\((k)\) R(A) is the radical of \(\mathrm{R}'(\mathrm{A})\) . (Use \((e)\) , \((f)\) , \((g)\) , \((h)\) and a Levi decomposition of A.) Let \(\mathrm{R}'(\mathrm{A}) = \mathrm{R}(\mathrm{A}) \oplus \mathrm{T}\) be a Levi decomposition of \(\mathrm{R}'(\mathrm{A})\) . By \((h)\) , \(\mathrm{R}'(\mathrm{A}) = \mathrm{R}(\mathrm{A}) \times \mathrm{T}\) . Then \(\mathrm{R}(\mathrm{T}) = \{0\}\) . Applying \((c)\) to the simple factors of T, conclude that \(\mathrm{T} = \{0\}\) . It has therefore been shown that \(\mathrm{R}(\mathrm{A}) = \mathrm{R}'(\mathrm{A})\) .

\begin{enumerate}
\def\labelenumi{(\alph{enumi})}
\setcounter{enumi}{11}
\item
  Let G be a connected real Lie group with Lie algebra A. Let \(\mathcal{S}(G)\) be the set of subgroups N of G such that there exists a decreasing sequence \((\mathrm{N}_m, \mathrm{N}_{m-1}, \ldots, \mathrm{N}_0)\) of subgroups with the following properties: \(\mathrm{N}_m = \mathrm{N}\) , \(\mathrm{N}_0 = \{\varepsilon\}\) , each \(\mathrm{N}_i\) is a normal connected Lie subgroup of G and \(\dim \mathrm{N}_i / \mathrm{N}_{i-1} = 1\) for all \(i > 0\) . Show that the integral subgroup R(G) of G with Lie algebra R(A) is the largest element of \(\mathcal{S}(G)\) . (Argue by induction on \(\dim \mathrm{R(A)}\) . Let I be a 1-dimensional ideal of A and N the corresponding integral subgroup of G. Pass to the quotient by N. If N is not closed, use Exercise 21 (a) of § 6.)
\item
  A Lie subgroup H of G is called extendable if there exists a Lie subgroup \(H_{1} \supset H\) such that \(\dim H_{1} = \dim H + 1\) . Let \(R'(G)\) denote the intersection of all the non-extendable connected Lie subgroups of G.
\end{enumerate}

Let B be a non-extendable subalgebra of A. Show that the corresponding integral subgroup of G is closed. (Use Proposition 5.) Deduce that

\[
\mathrm{L} (\mathrm{R} ^ {\prime} (\mathrm{G})) \subset \mathrm{R} (\mathrm{A}).
\]

whence \(\mathrm{R}^{\prime}(\mathrm{G})\subset\mathrm{R}(\mathrm{G})\) .

\begin{enumerate}
\def\labelenumi{(\alph{enumi})}
\setcounter{enumi}{13}
\tightlist
\item
  Show that \(\mathrm{R}(\mathrm{G}) = \mathrm{R}'(\mathrm{G})\) . (Reduce it to the case where \(\mathrm{R}'(\mathrm{G}) = \{e\}\) . Then use Exercise 22 of §6.)†
\end{enumerate}

\begin{enumerate}
\def\labelenumi{\arabic{enumi}.}
\setcounter{enumi}{28}
\tightlist
\item
  A real Lie group G is said to be of type (N) if it is finite-dimensional,
\end{enumerate}

nilpotent, connected and simply connected. If G is such a group and g is its Lie algebra, the mapping exp: g → G is an isomorphism when g is given the group structure defined by the Hausdorff law (cf.~Chapter II, §6, no. 5, Remark 3). Let log: G → g denote the inverse isomorphism.

\begin{enumerate}
\def\labelenumi{(\alph{enumi})}
\item
  Let V be a vector Q-subspace of g. Show the equivalence of the following conditions:
\item
  V is a Lie \(\mathbf{Q}\) -subalgebra of \(\mathfrak{g}\) ;
\end{enumerate}

\begin{enumerate}
\def\labelenumi{(\roman{enumi})}
\setcounter{enumi}{1}
\tightlist
\item
  \(\exp(V)\) is a subgroup of G.
\end{enumerate}

(Use Exercise 5 of § 6 of Chapter II.)

For a subgroup H of G to be of the form \(\exp(V)\) , where V is a vector Q-subspace of g, it is necessary and sufficient that H be saturated in G (Chapter II, §4, Exercise 14), i.e.~that the relations \(x \in G\) , \(x^{n} \in H\) , \(n \neq 0\) imply \(x \in H\) (loc. cit.). If this is so, show that H is an integral subgroup of G if and only if \(\log(H)\) is a Lie R-subalgebra of g.

\begin{enumerate}
\def\labelenumi{(\alph{enumi})}
\setcounter{enumi}{1}
\item
  Let V be a Lie Q-subalgebra of g of finite dimension m, let \((e_{1},\ldots,e_{m})\) be a basis of V over Q and let \(\Lambda\) be the subgroup of V generated by \(e_{1},\ldots,e_{m}\) . Using the fact that the Hausdorff law is polynomial, show that there exists an integer \(d\geqslant1\) such that, for every integer r which is a non-zero multiple of d, \(\exp(r\Lambda)\) is a subgroup of G. Let d be such an integer and r a multiple of d.~Show that \(\exp(r\Lambda)\) is discrete if and only if \((e_{1},\ldots,e_{m})\) is a free family over R, i.e.~if the canonical mapping of V \(\otimes_{q}R\) into g is injective. Suppose that this is the case; show that G/exp(r \(\Lambda\) ) is compact if and only if V \(\otimes_{q}R\to g\) is bijective, i.e.~if V is a Q-form of g. (To show that the condition is sufficient, argue by induction on the nilpotency class of g.)
\item
  Conversely, let \(\Gamma\) be a discrete subgroup of G, let \(\Gamma\) be its saturation in G (Chapter II, loc. cit.) and let \(\mathfrak{g}_{\Gamma} = \log (\Gamma) = \mathbf{Q}. \log (\Gamma)\) be the corresponding Lie \(\mathbf{Q}\) -subalgebra. Suppose that \(\mathbf{R}, \mathfrak{g} = \mathfrak{g}\) , i.e.~that \(\Gamma\) is contained in no distinct integral subgroup of G. Let \(z\) be an element \(\neq 1\) of the centre of \(\Gamma\) and let \(x = \log (z)\) . Show that \(x\) belongs to the centre of \(\mathfrak{g}\) . If \(\mathrm{X} = \exp (\mathbf{R}x)\) , show that \(\mathrm{X} / (\Gamma \cap \mathrm{X})\) is compact and that the image of \(\Gamma\) in \(\mathrm{G} / \mathrm{X}\) is discrete. Deduce, arguing by induction on dim G, that \(\mathrm{G} / \Gamma\) is compact and that \(\mathfrak{g}_{\Gamma}\) is a G-form of \(\mathfrak{g}\) . If \(\Lambda\) is a lattice of \(\mathfrak{g}_{\Gamma}\) , show that there exists an integer \(d \neq 0\) such that \(\Gamma\) contains \(\exp (d\Lambda)\) and that the index of \(\exp (d\Lambda)\) in \(\Gamma\) is then finite.
\end{enumerate}

Let H be an integral subgroup of G with Lie algebra \(\mathfrak{h}\) . Show that \(\mathrm{H} / (\mathrm{H} \cap \Gamma)\) is compact if and only if \(\mathfrak{h}\) is rational with respect to the \(\mathbf{Q}\) -structure \(g_{\Gamma}\) (in particular if \(\mathfrak{h}\) is one of the terms of the lower---or upper---central series of \(g\) ).

\begin{enumerate}
\def\labelenumi{(\alph{enumi})}
\setcounter{enumi}{3}
\item
  Let \(\Gamma\) be a discrete subgroup of \(G\) , let \(H\) be the smallest integral subgroup of \(G\) containing \(\Gamma\) and let \(\mathfrak{h}\) be its Lie algebra. Show that \(H / \Gamma\) is compact and that \(\mathfrak{g} \otimes_{\mathbf{q}} \mathbf{R} \to \mathfrak{g}\) is injective and has image \(\mathfrak{h}\) . (Apply (c) to the nilpotent group \(H\) .)
\item
  Let \(\Gamma\) be a discrete subgroup of \(G\) . Show that there exists a basis \((x_{1},\ldots ,x_{q})\) of \(L(G)\) with the following properties:
\item
  for all \(i \in [1, q]\) , \(\mathbf{R}x_i + \cdots + \mathbf{R}x_q\) is an ideal \(\mathfrak{n}_i\) of \(\mathbf{R}x_{i-1} + \cdots + \mathbf{R}x_q\) ;
\end{enumerate}

\begin{enumerate}
\def\labelenumi{(\roman{enumi})}
\setcounter{enumi}{1}
\tightlist
\item
  there exists \(p \in (1, q]\) such that \(\Gamma\) is the set of products
\end{enumerate}

\[
\exp (m _ {p} x _ {p}) \exp (m _ {p + 1} x _ {p + 1}) \dots \exp (m _ {q} x _ {q})
\]

where \(m_{p}, \ldots, m_{q}\) are in Z;

\begin{enumerate}
\def\labelenumi{(\roman{enumi})}
\setcounter{enumi}{2}
\tightlist
\item
  if \(G / \Gamma\) is compact, \(\{n_1, n_2, \ldots, n_q\}\) contains the lower central series of \(g\) . (Using (d) and Proposition 16, reduce it to the case where \(G / \Gamma\) is compact. Then argue by induction on \(\dim G\) , using (c).)
\end{enumerate}

\begin{enumerate}
\def\labelenumi{(\arabic{enumi})}
\setcounter{enumi}{29}
\tightlist
\item
  Give an example of a 7-dimensional real nilpotent Lie algebra with no basis with respect to which the constants of structure are rational (cf.~Chapter I, § 4, Exercise 18). Deduce that the corresponding group of type (N) (Exercise 29) has no discrete subgroup with compact subgroup. (Use Exercise 29.)
\end{enumerate}

\begin{enumerate}
\def\labelenumi{\arabic{enumi}.}
\setcounter{enumi}{30}
\item
  Let G and \(G'\) be two Lie groups of type (N) (Exercise 29) and let \(\Gamma\) be a discrete subgroup of G such that \(G/\Gamma\) is compact. Show that every homomorphism \(f: \Gamma \to G'\) can be extended uniquely to a Lie group morphism of G into \(G'\) (begin by extending f to the saturation \(\tilde{\Gamma}\) of \(\Gamma\) and derive a homomorphism of the Lie Q-algebra \(\log(\tilde{\Gamma})\) into the Lie algebra of \(G'\) , cf.~Exercise 29).
\item
  Let G be a Lie group of type (N) (Exercise 29) and let \(\Gamma\) be a discrete subgroup of G. Prove the equivalence of the following conditions:
\end{enumerate}

\begin{enumerate}
\def\labelenumi{(\alph{enumi})}
\item
  G/Γ is compact;
\item
  the measure of \(\mathrm{G} / \Gamma\) is finite (relative to a non-zero G-invariant positive measure);
\item
  every integral subgroup of G containing \(\Gamma\) is equal to G. (Use Exercise 29.)
\end{enumerate}

¶ 33. Let \(\Gamma\) be a group. Prove the equivalence of the following conditions: (a) \(\Gamma\) is nilpotent, torsion-free and finitely generated;

\begin{enumerate}
\def\labelenumi{(\alph{enumi})}
\setcounter{enumi}{1}
\item
  there exists a Lie group of type (N) (Exercise 29) containing \(\Gamma\) as a discrete subgroup;
\item
  there exists a Lie group G of type (N) (Exercise 29) containing \(\Gamma\) as a discrete subgroup and such that \(G / \Gamma\) is compact.
\end{enumerate}

(The equivalence of \(\langle b\rangle\) and \(\langle c\rangle\) follows from Exercise 29. The implication \((\varepsilon)\Rightarrow \langle a\rangle\) is proved by induction on \(\dim (\mathrm{G})\) by the method of Exercise 29 \((c)\) . To prove that \((a)\Rightarrow (c)\) , show first that the Lie \(\mathbf{Q}\) -algebra attached to the saturation of \(\Gamma\) (cf.~Chapter II, § 6, Exercise 4) is finite-dimensional; then take the tensor product of this algebra with \(\mathbf{R}\) and the corresponding group of type (N.)

\begin{enumerate}
\def\labelenumi{\arabic{enumi}.}
\setcounter{enumi}{33}
\tightlist
\item
  Let \(G\) be a finite-dimensional connected real or complex Lie group and \(\rho\) an analytic linear representation of \(G\) on a finite-dimensional complex vector space \(V\) . Suppose that \(\rho\) is semi-simple. The group \(G\) operates by automorphisms on the algebra \(S(V^*)\) of polynomial functions on \(V\) .
\end{enumerate}

\begin{enumerate}
\def\labelenumi{(\alph{enumi})}
\item
  Let \(\mathsf{S}(\mathbf{V}^{*})^{\mathbb{G}}\) be the set of G-invariant elements of \(\mathsf{S}(\mathbf{V}^{*})\) . There exists a projector \(p\) of \(S(V^{*})\) onto \(S(V^{*})^{\mathbb{G}}\) which commutes with the operations of \(G\) and which leaves stable every \(G\) -stable vector subspace of \(S(V^{*})\) .
\item
  Let I, J be ideals of \(S(V^*)\) which are stable under G. Let A (resp. B) be the set of zeros of I (resp. J) in V. Suppose that \(A \cap B = \varnothing\) . There then exists \(u \in I\) such that \(u\) is G-invariant and that \(u = 1\) in B. (By Commutative Algebra, Chapter 7, §3, no. 3, Proposition 2, there exists \(v \in I\) , \(w \in J\) such that \(v + w = 1\) Let \(u = pv\) . Show that \(u\) has the required properties.)
\end{enumerate}

\begin{enumerate}
\def\labelenumi{\arabic{enumi}.}
\setcounter{enumi}{34}
\tightlist
\item
  Let \(\mathbf{G}\) be a compact subgroup of \(\mathbf{GL}(n,\mathbf{R})\) .
\end{enumerate}

\begin{enumerate}
\def\labelenumi{(\alph{enumi})}
\item
  Let A, B be disjoint compact G-invariant subsets of \(\mathbf{R}^{\mathbf{n}}\) . There exists a G-invariant polynomial function \(u\) on \(\mathbf{R}^{\mathbf{n}}\) , such that \(u = 1\) on \(\mathrm{A}\) and \(u = 0\) on B. (There exists a continuous real-valued function \(v\) on \(\mathbf{R}^{\mathbf{n}}\) such that \(v = 1\) on A and \(v = 0\) on B. By the Stone-Weierstrass Theorem, there exists a polynomial function \(w\) on \(\mathbf{R}^{\mathbf{n}}\) such that \(|v - w| \leqslant \frac{1}{3}\) on \(\mathrm{A} \cup \mathrm{B}\) . Derive \(u\) by an integration process with respect to the normalized Haar measure of G.)
\item
  Let \(x \in \mathbf{R}^n\) . Then \(Gx\) is the set of zeros in \(\mathbf{R}^n\) of a finite number of G-invariant polynomials. (Use (a) and Corollary 2 to Theorem 1, no. 8.)
\end{enumerate}

\begin{enumerate}
\def\labelenumi{\arabic{enumi}.}
\setcounter{enumi}{35}
\tightlist
\item
  In \(\mathbf{SL}(2, \mathbf{R})\) , the elements which can be expressed in the form \((a, b)\) , where \(a, b\) in \(\mathbf{SL}(2, \mathbf{R})\) , are the elements \(\neq -1\) .
\end{enumerate}

¶ 37. Let G be a compact real Lie group and G' a finite-dimensional connected real Lie group. Suppose that L(G) is simple. Let \(\rho: G \to G'\) be a homomorphism of abstract groups. Suppose that there exists a neighbourhood V of \(e_{\alpha}\) in G such that \(\rho(V)\) is relatively compact. Then \(\rho\) is continuous. (Let V' be a neighbourhood of \(e_{\alpha'}\) in G'. Let V'\,' \(\subset\) V' be a neighbourhood of \(e_{\alpha'}\) , such that

\[
(x \in \mathrm{V} ^ {\prime \prime}, x _ {j} \in \rho (\mathrm{V}), y _ {j} \in \rho (\mathrm{V}) \text {   for   } 1 \leqslant j \leqslant n) \Rightarrow \left(\prod_ {i = 1} ^ {n} x _ {i} (y _ {i}, x) x _ {i} ^ {- 1} \in \mathrm{V} ^ {\prime}\right)
\]

where \(n = \dim G\) . It can be assumed that \(\rho\) is non-trivial. Then \(\operatorname{Ker} \rho\) is finite by Exercise 25. Therefore \(\rho(G)\) is not countable and hence not discrete. Hence there exists \(g \in G\) with non-open centralizer such that \(\rho(g) \in V''\) . Using Exercise 9 of §4 and its notation, \(\rho(M(g, n, V)) \subset V'\) and \(M(g, n, V)\) is a neighbourhood of \(e_G\) in G.)

\begin{enumerate}
\def\labelenumi{\arabic{enumi}.}
\setcounter{enumi}{37}
\tightlist
\item
  Let \(G\) be a finite-dimensional real Lie group and \(G_0\) its identity component. Consider the following property:
\end{enumerate}

\begin{enumerate}
\def\labelenumi{(\Alph{enumi})}
\setcounter{enumi}{5}
\tightlist
\item
  Ad(G) is closed in Aut L(G).
\end{enumerate}

Show that G has property (F) in each of the following cases:

\begin{enumerate}
\def\labelenumi{(\roman{enumi})}
\item
  G is connected and nilpotent;
\item
  every derivation of L(G) is inner;
\item
  \(\mathrm{G}_0\) has property (F) and \(\mathrm{G} / \mathrm{G}_0\) is finite;
\item
  G is an upper triangular group;
\item
  \(G_{0}\) is semi-simple.
\end{enumerate}

\begin{enumerate}
\def\labelenumi{\arabic{enumi}.}
\setcounter{enumi}{38}
\tightlist
\item
  Let G be a finite-dimensional connected real Lie Group, H an integral subgroup of G and \(\bar{H}\) its closure in G. Suppose that H has property (F) (Exercise 38).
\end{enumerate}

\begin{enumerate}
\def\labelenumi{(\alph{enumi})}
\item
  Let \(x \in \overline{\mathbf{H}}\) . Then \(\mathrm{Ad}_{\mathrm{L(G)}}x\) leaves \(\mathrm{L(H)}\) stable (no. 2, Proposition. 5); let \(u(x)\) be its restriction to \(\mathrm{L(H)}\) . Then \(u(\overline{\mathrm{H}}) = \mathrm{Ad}_{\mathrm{L(H)}}(\mathrm{H})\) . (Observe that \(\mathrm{Ad}_{\mathrm{L(H)}}(\mathrm{H})\) is dense in \(u(\overline{\mathrm{H}})\) and apply property (F).)
\item
  Let C be the centre of \(\overline{\mathrm{H}}\) . Then \(\overline{\mathrm{H}} = \mathrm{C.H}\) and C is the closure of the centre of H. (By (a), if \(x \in \overline{\mathrm{H}}\) , there exists \(y \in \mathrm{H}\) such that \(\mathrm{Ad}_{\mathrm{L(H)}} x = \mathrm{Ad}_{\mathrm{L(H)}} y\) . Then \(\mathrm{Ad}(x^{-1}y)\) is the identity on L(H), hence \(x^{-1}y \in Z_{\overline{\mathrm{H}}}(\mathrm{H})\) and \(x \in Z_{\overline{\mathrm{H}}}(\mathrm{H})\) . H. By continuity, \(Z_{\overline{\mathrm{H}}}(\mathrm{H})\) is equal to C. The group C is closed in \(\overline{\mathrm{H}}\) and hence in G; therefore \(\overline{\mathrm{C}} \cap \overline{\mathrm{H}} \subset \mathrm{C}\) . Let \(x \in \mathrm{C}\) . There exists a sequence \((x_n)\) in H such that \(x_n\) tends to \(x\) . Then \(\mathrm{Ad}_{\mathrm{L(H)}} x_n\) tends to 1 and hence there exists a sequence \((y_n)\) in H such that \(y_n\) tends to \(e\) and \(\mathrm{Ad}_{\mathrm{L(H)}} y_n = \mathrm{Ad}_{\mathrm{L(H)}} x_n\) . Then \(y_n^{-1}x_n \in \mathrm{C} \cap \mathrm{H}\) and \(y_n^{-1}x_n\) tends to \(x\) .)
\item
  \(\mathbf{H} = \overline{\mathbf{H}}\) if and only if the centre of \(\mathbf{H}\) is closed in \(\mathbf{G}\) . (Use (b).)
\end{enumerate}

\begin{enumerate}
\def\labelenumi{\arabic{enumi}.}
\setcounter{enumi}{39}
\tightlist
\item
  Let \(\mathbf{H}\) be a finite-dimensional real Lie group and \(\mathbf{H}_0\) its identity component.
\end{enumerate}

\begin{enumerate}
\def\labelenumi{(\alph{enumi})}
\item
  Suppose that \(H_{0}\) satisfies condition (F) of Exercise 38, that \(H/H_{0}\) is finite and that the centre of \(H_{0}\) is compact. Let G be a finite-dimensional real Lie group and \(f:H\to G\) a continuous homomorphism with discrete kernel. Then \(f(\mathrm{H})\) is closed in G. (Reduce it to the case where H is connected. As the kernel of f is discrete, the centre of \(f(\mathrm{H})\) is the image under f of the centre of H (Lemma 1) and hence is closed in G. Apply Exercise 39 (c).)
\item
  Suppose that \(H_{0}\) is semi-simple and that \(H/H_{0}\) is finite. Then the image of H under a finite-dimensional continuous linear representation is closed in the linear group. (Apply (a) and Exercise 7 (b).)
\end{enumerate}

\begin{enumerate}
\def\labelenumi{\arabic{enumi}.}
\setcounter{enumi}{40}
\tightlist
\item
  Let \(G\) be a finite-dimensional connected real Lie group and \(\rho: G \to \mathbf{GL}(n, \mathbf{C})\) a continuous homomorphism with finite kernel. Then \(\rho((G, G))\) is closed in \(\mathbf{GL}(n, \mathbf{C})\) . (Using Exercises 7 (b) and 40 (a), reduce it to the case where \(G\) is the semi-direct product of a semi-simple group \(S\) and its radical \(R\) . Every element of \(\rho((G, R))\) is unipotent and hence \(\rho((G, R))\) is closed. The vector subspace \(V_1\) of fixed points of \(\rho((G, R))\) is \(\neq \{0\}\) . It is stable under \(\rho(G)\) . By considering \(\mathbf{C}^n / V_1\) arguing by induction on \(n\) and using the complete reducibility of \(\rho(S)\) , it is seen that it is possible to write
\end{enumerate}

\[
\mathbf {C} ^ {n} = \mathrm{V} _ {1} \oplus \dots \oplus \mathrm{V} _ {q}
\]

where each \(V_{i}\) is stable under \(\rho(S)\) and \(\rho((G,R))(V_{i})\subset V_{1}\oplus\cdots\oplus V_{i-1}\) for all i. On the other hand, \(\rho(S)\) is closed (Exercise 40 (b)). Finally, \(\rho((G,G))=\rho(S)\cdot\rho((G,R)).\)

\begin{enumerate}
\def\labelenumi{\arabic{enumi}.}
\setcounter{enumi}{41}
\tightlist
\item
  Let \(G\) be a finite-dimensional connected real Lie group. Suppose that
\end{enumerate}

G admits an injective continuous linear representation \(\rho: \mathrm{G} \to \mathbf{GL}(n, \mathbf{C})\) . Then G admits a linear representation \(\sigma: \mathrm{G} \to \mathbf{GL}(n', \mathbf{C})\) which is a homomorphism of G onto a closed subgroup of \(\mathbf{GL}(n', \mathbf{C})\) . (By Exercise 41, \(\rho((\mathrm{G}, \mathrm{G}))\) is closed in \(\mathbf{GL}(n, \mathbf{C})\) and hence (G, G) is closed in G. Let \(\rho: \mathrm{G} \to (\mathrm{G}, \mathrm{G})\) be the canonical morphism and \(\tau\) an injective linear representation with closed image \(\mathrm{G}/(\mathrm{G}, \mathrm{G})\) which is connected and commutative. Take \(\sigma = \rho \oplus (\tau \circ p)\) .)

\exercisesfor{10}

\begin{enumerate}
\def\labelenumi{\arabic{enumi}.}
\item
  Let \(G\) be a finite-dimensional connected real Lie group. Show that the canonical mapping of \(\operatorname{Aut}(G)\) into \(\operatorname{Aut}(L(G))\) is not in general surjective (take \(G = T\) ).
\item
  Suppose that \(\mathbf{K} = \mathbf{Q}_p\) . Let \(\mathbf{G}\) be a finite-dimensional Lie group. Show that the following conditions are equivalent:
\end{enumerate}

\begin{enumerate}
\def\labelenumi{(\alph{enumi})}
\item
  there exist \(x_{1}, x_{2}, \ldots, x_{n}\) in \(G\) such that the subgroup of \(G\) generated by \(\{x_{1}, \ldots, x_{n}\}\) is dense in \(G\) ;
\item
  G is generated by a compact subset.
\end{enumerate}

(To prove that (b) implies (a), observe that, if \((e_1, \ldots, e_n)\) is a basis of \(\mathrm{L}(\mathrm{G})\) , \((\exp \mathbf{Z}_p e_1)(\exp \mathbf{Z}_p e_2) \ldots (\exp \mathbf{Z}_p e_n)\) is a neighbourhood of \(e\) .)

\begin{enumerate}
\def\labelenumi{\arabic{enumi}.}
\setcounter{enumi}{2}
\tightlist
\item
  Let \(G\) be the set of \((x, y) \in \mathbf{Q}_p \times \mathbf{Q}_p\) such that \(|y| \leqslant 1\) . This is an open subgroup of \(\mathbf{Q}_p \times \mathbf{Q}_p\) . Then \(L(G) = \mathbf{Q}_p \times \mathbf{Q}_p\) . Let \(\alpha\) be the infinitesimal automorphism \((x, y) \mapsto (0, x)\) . Then the conclusion of Proposition 3 is false.
\end{enumerate}

\begin{enumerate}
\def\labelenumi{(\arabic{enumi})}
\setcounter{enumi}{3}
\tightlist
\item
  Let \(\mathbf{K}\) be a quadratic extension of \(\mathbf{Q}_p\) and let \(\omega \in \mathbf{K} - \mathbf{Q}_p\) . Let \(\mathrm{G} = \mathbf{Q}_p + \mathbf{Z}_p\omega\) , considered as a Lie group over \(\mathbf{K}\) . The only automorphisms of \(\mathrm{G}\) are the \((x,y)\mapsto (\lambda x,y)\) where \(\lambda\) is invertible in \(\mathbf{Z}_p\) . The set of these automorphisms cannot be given a Lie group structure over \(\mathbf{K}\) with the properties of Theorem 1.
\end{enumerate}

\specialchapter{Historical note (Chapters I to III)}

\hnsection{GENESIS}

The theory, called for nearly a century ``theory of Lie groups'', was essentially developed by one mathematician: Sophus Lie.

Before embarking on the history, we summarize briefly various earlier research which prepared the way.

\textbf{(a) Transformation groups (Klein-Lie, 1869--1872)}

About 1860 the theory of permutation groups of a finite set was developing and beginning to be used (Serret, Kronecker, Mathieu, Jordan). On the other hand, the theory of invariants, which was in full flight, was familiarizing mathematicians with some infinite sets of geometric transformations stable under composition (notably linear or projective transformations). But before the work of Jordan {[}7{]} in 1868 on ``groups of movements'' (closed subgroups of the group of displacements of 3-dimensional Euclidean space), no conscious link seems to have been established between these two currents of ideas.

In 1869, the young Felix Klein (1849--1925), a pupil of Plücker, established a friendship in Berlin with the Norwegian Sophus Lie (1842--1899), who was a few years older, brought together by their common interest in Plücker's ``line geometry'' and in particular the theory of line complexes. It was about this time that Lie conceived one of his most original ideas, the introduction of the notion of invariant into Analysis and Differential Geometry; one of the sources was his observation that the classical methods of integration ``by quadratures'' of differential equations depended entirely on the fact that the equation is invariant under a ``continuous'' family of transformations. 1869 is the date of the first work (edited by Klein) where Lie used this idea; there he studied the ``Reye complex'' (set of four lines cutting the faces of a tetrahedron in four points with a given cross-ratio) and the curves and surfaces which admit as tangents lines of this complex {[}3a{]}: his method depends on the invariance of the Reye complex under the 3-parameter commutative group (maximal torus of PGL(4, C)) which leaves the vertices of the tetrahedron invariant. The same idea dominated the work Klein and Lie wrote together when they were in Paris in the spring of 1870 {[}1, a{]}; there they essentially determined the commutative connected subgroups of the projective group of the plane PGL(3, C) and studied the geometric properties of their orbits (under the name of curves or surfaces V); this gave them, by a uniform procedure, properties of various curves, algebraic or transcendental, such as \(y = cx^m\) or the logarithmic spirals. Both works served to underline the profound impression made on them by the theories of Galois and Jordan (Jordan's commentary on Galois had appeared in Math. Annalen in 1869; on the other hand, Lie had heard talk of Galois theory as early as 1863). Klein, who in 1871 became interested in non-Euclidean geometries, saw there the start of his research on a classification principle for all known geometries which was to lead him in 1872 to the ``Erlangen Programme''. For his part, Lie, in a letter of 1873 to A. Mayer ({[}3{]}, vol.~V, p.~584), dated his ideas on transformation groups from his sojourn in Paris and in a work of 1871 ({[}3 b{]} p.~208) he was already using the term ``transformation group'' and explicitly posed the problem of the determination of all subgroups (``continuous or discontinuous'') of GL(n, C). To be truthful, Klein and Lie had both experienced some difficulty in entering this new mathematical universe and Klein spoke of Jordan's newly appeared ``Treatise'' as a ``book sealed with seven seals'' ({[}2{]}, p.~51); he wrote moreover concerning {[}1 a{]} and {[}1 b{]}: ``To Lie belongs all the credit for the heuristic idea of a continuous group of operators, in particular everything concerning integration of differential equations and partial derivatives. All the notions he developed later in his theory of continuous groups were already there in embryo, but were however so little elaborated, that only after long conversations could I convince him of many details, for example to begin with the very existence of the curves V'' ({[}2{]}, p.~415).

\textbf{(b) Infinitesimal transformations}

The concept of an ``infinitely small'' transformation goes back at least to the beginnings of Infinitesimal Calculus; we know that Descartes discovered the instantaneous centre of rotation when admitting that ``in the infinitely small'' every plane movement can be likened to a rotation: the elaboration of Analytical Mechanics in the 18th century is entirely founded on similar ideas. In 1851, Sylvester, seeking to form invariants of the linear group \(\mathbf{GL}(3,\mathbf{C})\) and some of its subgroups, gave the parameters appearing in his matrices ``infinitely small'' increases of the forms \(\alpha_{j}dt\) and expressed the fact that a function \(f((z_{j}))\) was invariant by writing the equation \(f((z_{j}+\alpha_{j}(dt))=f((z_{j}))\) ; this gave him for f the linear partial differential equation Xf=0, where

\[
\mathrm{X} f = \sum_ {j} \alpha_ {j} \frac {\partial f}{\partial z _ {j}},\tag{1}
\]

where X is therefore a differential operator, the ``derivative in the direction of the direction parameters \(\alpha_{j}\) ({[}5{]}, vol.~3, pp.~326 and 327); Sylvester seemed to think that here was a general principle of considerable importance but appears never to have returned to the question. A little later, Cayley ({[}6{]}, vol.~II, pp.~164--178) proceeded similarly for the invariants of \(\mathbf{SL}(2,\mathbf{C})\) under certain representations of this group and showed that they are the solutions of two first order partial differential equations \(Xf=0\) , \(Yf=0\) , where X and Y are obtained as above from ``infinitely small'' transformations

\[
\left( \begin{array}{c c} 0 & 0 \\ d t & 0 \end{array} \right) \quad \text { and } \quad \left( \begin{array}{c c} 0 & d t \\ 0 & 0 \end{array} \right).
\]

In modern terms, this is expressed by the fact that X and Y generate the Lie algebra \(\mathfrak{sl}(2,\mathbf{C})\) ; moreover Cayley calculated the bracket XY -- YX explicitly and showed that it was also derived from an ``infinitely small'' transformation.

In his memoir of 1868 on groups of movements {[}7{]}, Jordan used from beginning to end the concept of ``infinitely small transformation'', but exclusively from a geometric point of view. He is no doubt responsible for the idea of a one-parameter group ``generated'' by an infinitely small transformation: for Jordan, it was the set of transformations obtained by ``suitably repeating'' the infinitely small transformation (loc. cit., p.~243). Klein and Lie, in their memoir, used the same expression ``repeated infinitely small transformation'' {[}1 b{]}, but the context shows that they understood by that an integration of a differential system. If the one-parameter group they considered consists of transformations \(x' = f(x, y, t)\) , \(y' = g(x, y, t)\) , the corresponding ``infinitely small transformation'' is given by

\[
d x = p (x, y) d t, \quad d y = q (x, y) d t
\]

where \(p(x,y) = \frac{\partial f}{\partial t} (x,y,t_0), q(x,y) = \frac{\partial g}{\partial t} (x,y,t_0)\) and \(t_0\) corresponds to the

identity transformation of the group. As Klein and Lie knew the functions f and g explicitly, they had no difficulty in verifying that the functions

\[
t \mapsto f (x, y, t) \quad \text { and } \quad t \mapsto g (x, y, t)
\]

give in parametric form the integral curve of the differential equation

\[
q (\xi , \eta) d \xi = p (\xi , \eta) d \eta
\]

passing through the point \((x,y)\) , but they give no general argument; moreover they nowhere use this fact in the remainder of their memoir.

\textbf{(c) Contact transformations}

In the next two years, Lie seemed to abandon the theory of transformation groups (although he remained in very close contact with Klein, who published his ``Programme'' in 1872) to study contact transformations, integration of first order partial differential equations and the relations between these two theories. We are not concerned here with the history of these questions and we shall confine ourselves to mentioning a few points which seem to have played an important role in the genesis of the theory of transformation groups.

The notion of contact transformation generalizes both point transformations and inverse polar transformations. Roughly a contact transformation† on \(\mathbf{C}^{\mathrm{n}}\) is an isomorphism of an open subset \(\Omega\) of the manifold \(T^{\prime}(\mathbf{C}^{\mathrm{n}})\) of cotangent vectors to \(\mathbf{C}^{\mathrm{n}}\) onto another open subset \(\Omega^{\prime}\) of \(T^{\prime}(\mathbf{C}^{\mathrm{n}})\) mapping the canonical 1-form of \(\Omega\) to that of \(\Omega^{\prime}\) . In other words, if \((x_{1},\ldots ,x_{n},p_{1},\ldots ,p_{n})\) denote the canonical coordinates of \(T^{\prime}(\mathbf{C}^{\mathrm{n}})\) , a contact transformation is an isomorphism \((x_{1},p_{1})\mapsto (\mathbf{X}_{\mathrm{i}},\mathbf{P}_{\mathrm{i}})\) satisfying the relation \(\sum_{i = 1}^{n}\mathbf{P}_{i}d\mathbf{X}_{i} = \sum_{i = 1}^{n}p_{i}dx_{i}\) . Such transformations occur in the study of the integration of partial differential equations of the form

\[
\mathrm{F} \left(x _ {1}, x _ {2}, \dots , x _ {n}, \frac {\partial z}{\partial x _ {1}}, \dots , \frac {\partial z}{\partial x _ {n}}\right) = 0.\tag{2}
\]

Lie became familiar during his research on these questions with the manipulation of Poisson brackets

\[
(f, g) = \sum_ {i = 1} ^ {n} \left(\frac {\partial f}{\partial x _ {i}} \frac {\partial g}{\partial p _ {i}} - \frac {\partial g}{\partial x _ {i}} \frac {\partial f}{\partial p _ {i}}\right)\tag{3}
\]

and the brackets \(\ddagger_{\xi}\) {[}X, Y{]} = XY - YX of differential operators of type (1); he interprets the Poisson bracket (3) as the effect on f of a transformation of type (1) associated with g and observes on this occasion that the Jacobi identity for Poisson brackets means that the bracket of the differential operators corresponding to f and g is associated with the bracket (g, h). Research into functions g such that (F, g) = 0, which occurs in Jacobi's method for integrating the partial differential equation (2), became for Lie the study of infinitesimal contact transformations which leave the given equation invariant. Finally, Lie was led to study sets of functions \((u_{j})_{1 \in j \leq m}\) of the \(x_{i}\) and \(p_{i}\) such that the brackets \((u_{j}, u_{k})\) are functions of the \(u_{h}\) and called these sets ``groups'' (they had already essentially been studied by Jacobi) {[}3 c{]}.

\hnsection{CONTINUOUS GROUPS AND INFINITESIMAL TRANSFORMATIONS}

Suddenly, in the autumn of 1873, Lie again took up the study of transformation groups and obtained decisive results. Insofar as it is possible to follow his line of thought from a few letters to A. Mayer written in 1873--1874 ({[}3{]}, vol.~5, pp.~584--608), he started from a ``continuous group'' of transformations on n variables

\[
x _ {i} ^ {\prime} = f _ {i} (x _ {1}, \dots , x _ {n}, a _ {1}, \dots , a _ {r}) \quad (1 \leqslant i \leqslant n)\tag{4}
\]

depending effectively† on r parameters \(a_{1},\ldots,a_{r}\) ; he observed that, if the transformation (4) is the identity for the values \(a_{1}^{0},\ldots,a_{r}^{0}\) of the parameters,‡ then the first order Taylor expansions of the \(x_{i}\) :

\[
\begin{array}{r l} f _ {i} (x _ {1}, \dots , x _ {n}, a _ {1} ^ {0} + z _ {1}, \dots , a _ {r} ^ {0} + z _ {r}) & \\ = x _ {i} + \sum_ {k = 1} ^ {r} z _ {k} X _ {k i} (x _ {1}, \dots , x _ {n}) + \dots & (1 \leqslant i \leqslant n) \end{array} \tag {5}
\]

give a ``generic'' infinitely small transformation depending linearly on the r parameters \(z_{j}\)

\[
d x _ {i} = \left(\sum_ {k = 1} ^ {r} z _ {k} \mathrm{X} _ {k i} (x _ {1}, \dots , x _ {n})\right) d t \quad (1 \leqslant i \leqslant n).\tag{6}
\]

Proceeding as in his memoir with Klein, Lie integrated the differential system

\[
\frac {d \xi_ {1}}{\sum_ {k} z _ {k} X _ {k 1} (\xi_ {1} , \dots , \xi_ {n})} = \dots = \frac {d \xi_ {n}}{\sum_ {k} z _ {k} X _ {k n} (\xi_ {1} , \dots , \xi_ {n})} = d t,\tag{7}
\]

which gave him, for each point \((z_{1},\ldots,z_{r})\) a one-parameter group

\[
t \mapsto x _ {i} ^ {\prime} = g _ {i} (x _ {1}, \dots , x _ {n}, z _ {1}, \dots , z _ {r}, t) \quad (1 \leqslant i \leqslant n)\tag{8}
\]

such that \(g_{i}(x_{1},\ldots,x_{n},z_{1},\ldots,z_{r},0)=x_{i}\) for all i. He showed by an ingenious method, using the fact that the transformations (4) form a set which is stable under composition, that the one-parameter group (8) is a subgroup of the given group {[}3 d{]}. The new idea, the key to the whole theory, is to take the Taylor expansions of the functions (4) to the second order. The progress of his argument was quite confused and heuristic ({[}3 d{]} and {[}3{]}, vol.~5, pp.~600--601); it can be presented as follows. For sufficiently small \(z_{r}\) , let t=1 in (8); thus new parameters \(z_{1},\ldots,z_{r}\) are obtained for the transformations of the group (this is in fact the first appearance of the ``canonical parameters''). Then by definition, using (7),

\[
\frac {\partial g _ {i}}{\partial t} = \sum_ {k} z _ {k} \mathrm{X} _ {k i} (x _ {1} ^ {\prime}, \dots , x _ {n} ^ {\prime})
\]

whence

\[
\begin{array}{r l} \frac {\partial^ {2} g _ {i}}{\partial t ^ {2}} & = \sum_ {k, j} z _ {k} \frac {\partial \mathrm{X} _ {k i}}{\partial x _ {j}} (x _ {1} ^ {\prime}, \ldots , x _ {n} ^ {\prime}) \frac {\partial x _ {j} ^ {\prime}}{\partial t} \\ & = \sum_ {k, j} z _ {k} \frac {\partial \mathrm{X} _ {k i}}{\partial x _ {j}} (x _ {1} ^ {\prime}, \ldots , x _ {n} ^ {\prime}) \left(\sum_ {n} z _ {h} \mathrm{X} _ {h j} (x _ {1} ^ {\prime}, \ldots , x _ {n} ^ {\prime})\right) \end{array}
\]

which gives

\[
\begin{array}{r} x _ {i} ^ {\prime} = x _ {i} + \Big (\sum_ {k} z _ {k} \mathrm{X} _ {k i} (x _ {1}, \ldots , x _ {n}) \Big) t \\ + \frac {1}{2} \Big (\sum_ {k, n, j} z _ {k} z _ {h} \frac {\partial \mathrm{X} _ {k i}}{\partial x _ {j}} (x _ {1}, \ldots , x _ {n}) \mathrm{X} _ {h j} (x _ {1}, \ldots , x _ {n}) \Big) t ^ {2} + \dots , \end{array}
\]

whence, for \(t = 1\) , the Taylor expansions in the parameters \(z_{j}\)

\[
x _ {i} ^ {\prime} = x _ {i} + \left(\sum_ {j} z _ {k} \mathrm{X} _ {k i}\right) + \frac {1}{2} \left(\sum_ {k, h, j} z _ {k} z _ {h} \mathrm{X} _ {h j} \frac {\partial \mathrm{X} _ {k i}}{\partial x _ {j}}\right) + \dots (1 \leqslant i \leqslant n).\tag{9}
\]

We abridge these relations to \(x' = \mathbf{G}(x, z)\) relating the vectors

\[
x = (x _ {1}, \dots , x _ {n}), \quad x ^ {\prime} = (x _ {1} ^ {\prime}, \dots , x _ {n} ^ {\prime}), \quad z = (z _ {1}, \dots , z _ {r});
\]

the fundamental stability property of the set of these transformations under composition can be written

\[
\mathrm{G} (\mathrm{G} (x, u), v) = \mathrm{G} (x, \mathrm{H} (u, v))\tag{10}
\]

where \(\mathbf{H} = (\mathbf{H}_1, \ldots, \mathbf{H}_r)\) is independent of \(x\) ; it is immediate that \(\mathrm{H}(u, 0) = u\) and \(\mathrm{H}(0, v) = v\) , whence the expansions

\[
\mathrm{H} _ {i} (u, v) = u _ {i} + v _ {i} + \frac {1}{2} \sum_ {h, k} c _ {i k h} u _ {h} v _ {k} + \dots ,\tag{11}
\]

where the terms omitted are not linear in u or v. Transforming (10), using (9) and (11) and then comparing the terms in \(u_{h}v_{k}\) of the two sides. Lie obtained the relations

\[
\sum_ {j = 1} ^ {n} \left(\mathrm{X} _ {h j} \frac {\partial \mathrm{X} _ {k i}}{\partial x _ {j}} - \mathrm{X} _ {k j} \frac {\partial \mathrm{X} _ {h i}}{\partial x _ {j}}\right) = \sum_ {l = 1} ^ {r} c _ {l h k} \mathrm{X} _ {l i} (1 \leqslant h, k \leqslant r, 1 \leqslant i \leqslant n). \tag{12}
\]

His experience with the theory of partial differential equations led him to write these conditions in a simpler form: following the pattern of (1), he associated with each of the infinitely small transformations obtained by setting \(z_{k} = 1\) , \(z_{h} = 0\) for \(h \neq k\) in (6), the differential operator

\[
\mathrm{A} _ {k} (f) = \sum_ {i = 1} ^ {n} \mathrm{X} _ {k i} \frac {\partial f}{\partial x _ {i}},\tag{13}
\]

and rewrote conditions (12) in the form

\[
[ \mathrm{A} _ {h}, \mathrm{A} _ {k} ] = \sum_ {l} c _ {l h k} \mathrm{A} _ {l},\tag{14}
\]

the corner stone of his theory. Until then he used the terms ``infinitely small transformation'' and ``infinitesimal transformation'' indifferently (e.g.~{[}3 e{]}); the simplicity of relations (14) led him to call the operator (13) the ``symbol'' of the infinitesimal transformation \(dx_{i} = X_{ki}dt\) ( \(1 \leqslant i \leqslant n\) ) {[}3 d{]} and very soon he called the operator (13) itself the ``infinitesimal transformation'' ({[}3 d{]} and {[}3{]}, vol.~5, p.~589).

He then became aware of the close links which united the theory of ``continuous groups'' with his earlier research on contact transformations and partial differential equations. This bringing together filled him with enthusiasm: ``My earlier works were as it were all ready there waiting to found the new theory of transformation groups'' he wrote to Mayer in 1874 ({[}3{]}, vol.~5, p.~586).

In the following years, Lie pursued the study of transformation groups. Besides the general theorems summarized below (§ III), he obtained some more special results: the determination of the transformation groups of the line and plane, the subgroups of low codimension in the projective groups, the groups in at most 6 parameters, etc. He did not abandon differential equations for long. In fact, it seems that, for him, the theory of transformation groups was a tool for integrating differential equations, where the transformation group played a role analogous to that of the Galois group of an algebraic equation.†

We note that this research led him equally to the introduction of certain transformation sets in an infinity of parameters which he called ``infinite continuous groups'' \(^{‡}\) ; he reserved the name ``finite continuous groups'' for transformation groups in a finite number of parameters of type (4) above.

\hnsection{THE LIE GROUPS-LIE ALGEBRAS ``DICTIONARY''}

The theory of ``finite continuous'' groups, developed by Lie in numerous memoirs beginning in 1874, is expounded systematically in the imposing treatise ``Theorie der Transformationsgruppen'' ({[}4{]}, 1888--1893{]}), written in collaboration with F. Engel§; it is the object of study in the first volume and the last five chapters of the third, the second being devoted to contact transformations.

As the title indicates, this work is only concerned with transformation groups in the sense of equations (4), where the space of ``variables'' \(x_{i}\) and the space of ``parameters'' \(a_{i}\) play initially such important roles. On the other hand, the concept of ``abstract'' group had not been clearly isolated at that period; when in 1885 ({[}3/j{]}, §5) Lie noted that in the notation of (10) the equation \(w = \mathrm{H}(u, v)\) which gives the parameters of the composition of two transformations of the group defines a new group, he considered it as a transformation group on the space of parameters, thus obtaining what he called the ``parameter group'' (he even obtained two, which are just the group of left translations and the group of right translations†).

The variables \(x_{i}\) and the parameters \(a_{i}\) in equations (4) were in principle assumed to be complex (except in Chapters XIX--XXIV of volume 3) and the functions \(f_{i}\) to be analytic; Lie and Engel were of course aware of the fact that these functions are not in general defined for all complex values of the \(x_{i}\) and \(a_{i}\) and that therefore the composition of such functions raises serious difficulties ({[}4{]}, vol.~1, pp.~15--17, pp.~33--40 and passim); and although throughout they almost always wrote as if the composition of the transformations they were studying was possible without restriction, this was no doubt for the convenience of the results and they explicitly reaffirmed the ``local'' point of view whenever necessary (cf.~loc. cit., p.~168 or 189 for example or ibid., vol.~3, p.~2, note at the bottom of the page); in other words, the mathematical object they studied is close to what we call in this treatise a law chunk of operation. They did not refrain, on occasions, from considering global groups, for example the 4 series of classical groups ({[}4{]}, vol.~3, p.~682), but do not appear to have asked themselves the question of what in general constitutes a ``global group''; they were content to obtain, for the ``parameters'' of the classical groups (the ``variables'' of these groups introduced no difficulty, since the transformations in question are linear transformations of \(\mathbf{C}^n\) ) ``local'' systems of parameters in a neighbourhood of the identity transformation, without worrying about the domain of validity of the formulae they were writing down. They however set themselves a problem which arises neatly out of the local theory: the study of ``mixed'' groups, that is groups with a finite number of connected components, such as the orthogonal group ({[}4{]}, vol.~1, p.~7). They presented this study as a study of a set of transformations stable under composition and passage to the inverse which is the union of sets \(\mathrm{H}_j\) each of which is described by systems of functions ( \(f_i^{(D)}\) ) as in (4); the number of (essential) parameters of each \(\mathrm{H}_j\) is even a priori assumed to depend on \(j\) , but they showed that in fact this number is the same for all the \(\mathrm{H}_j\) . Their principal result was then the existence of a finite continuous group G such that \(\mathrm{H}_j = \mathrm{G} \circ h_j\) for some \(h_j \in \mathrm{H}_j\) and for all \(j\) ; they also established that G is normal in the mixed group and noted that the determination of the invariants of the latter reduces to that of the invariants of G and a discontinuous group ({[}4{]}, vol.~1, Chapter 18).

The general theory developed in {[}4{]} ended (without the authors saying so very systematically) by achieving a ``dictionary'' translating the properties of ``finite continuous'' groups into those of the set of their infinitesimal transformations. It is based on the ``three theorems of Lie'', each of which consists of an assertion and its converse.

The first theorem ({[}4{]}, vol.~1, pp.~33 and 72 and vol.~3, p.~563) affirms in the first place that if in (4) the parameters are effective, the functions \(f_{i}\) satisfy a system of partial differential equations of the form

\[
\frac {\partial f _ {i}}{\partial a _ {j}} = \sum_ {k = 1} ^ {r} \xi_ {k i} (f (x, a)) \psi_ {k j} (a) \quad (1 \leqslant i \leqslant n)\tag{15}
\]

where the matrix \((\xi_{kl})\) is of maximum rank and \(\det(\psi_{kj}) \neq 0\) ; conversely, if the functions \(f_{i}\) have this property, formulae (4) define a group germ of transformations.

The second theorem ({[}4{]}, vol.~1, pp.~149 and 158, and vol.~3, p.~590) gives relations between the \(\xi_{ki}\) on the one hand and the \(\psi_{ij}\) on the other: the conditions on the \(\xi_{ki}\) can be written in the form

\[
\sum_ {k = 1} ^ {n} \left(\xi_ {i k} \frac {\partial \xi_ {f l}}{\partial x _ {k}} - \xi_ {j k} \frac {\partial \xi_ {t l}}{\partial x _ {k}}\right) = \sum_ {k = 1} ^ {r} c _ {i j} ^ {k} \xi_ {k l} (1 \leqslant i, j \leqslant r, 1 \leqslant l \leqslant n)\tag{16}
\]

where the \(c_{ii}^{k}\) are constants \((1\leqslant i,j,k\leqslant r)\) skew-symmetric in \(i,j\) . The conditions on the \(\psi_{ij}\) , in the form given by Maurer {[}10{]}, are:

\[
\frac {\partial \psi_ {k l}}{\partial a _ {m}} - \frac {\partial \psi_ {k m}}{\partial a _ {l}} = \frac {1}{2} \sum_ {1 \leqslant i, j \leqslant r} c _ {i j} ^ {k} (\psi_ {i l} \psi_ {j m} - \psi_ {j l} \psi_ {i m}) (1 \leqslant k, l, m \leqslant r).\tag{17}
\]

By introducing the contragradient matrix \((\alpha_{ij})\) of \((\psi_{ij})\) and the infinitesimal transformations

\[
\mathbf {X} _ {k} = \sum_ {i = 1} ^ {n} \xi_ {k i} \frac {\partial}{\partial x _ {i}}, \quad \mathbf {A} _ {k} = \sum_ {j = 1} ^ {r} \alpha_ {k j} \frac {\partial}{\partial a _ {j}} (1 \leqslant k \leqslant r)\tag{18}
\]

\begin{enumerate}
\def\labelenumi{(\arabic{enumi})}
\setcounter{enumi}{15}
\item
  and (17) can be written respectively:
\item
\end{enumerate}

\[
[ \mathrm{X} _ {i}, \mathrm{X} _ {j} ] = \sum_ {k = 1} ^ {r} c _ {i j} ^ {k} \mathrm{X} _ {k} (1 \leqslant i, j \leqslant r).\tag{20}
\]

\[
[ \mathbf {A} _ {i}, \mathbf {A} _ {j} ] = \sum_ {k = 1} ^ {r} c _ {i j} ^ {k} \mathbf {A} _ {k}
\]

Conversely, if r infinitesimal transformations \(X_{k}\) ( \(1 \leqslant k \leqslant r\) ) are given which are linearly independent and satisfy conditions (19), the one-parameter subgroups generated by these transformations generate a transformation group in r essential parameters.

Finally, the third theorem ({[}4{]}, vol.~1, pp.~170 and 297 and vol.~3, p.~597) reduces the determination of the systems of infinitesimal transformations \((\mathbf{X}_k)_{1 \leqslant k \leqslant r}\) satisfying (19) to a purely algebraic problem: the following must hold:

\begin{enumerate}
\def\labelenumi{(\arabic{enumi})}
\setcounter{enumi}{20}
\tightlist
\item
\end{enumerate}

\[
c _ {i j} ^ {k} + c _ {j i} ^ {k} = 0\tag{22}
\]

\[
\sum_ {l = 1} ^ {r} \left(c _ {i l} ^ {m} c _ {j k} ^ {l} + c _ {k i} ^ {m} c _ {i j} ^ {l} + c _ {j l} ^ {m} c _ {k l} ^ {l}\right) = 0 \quad (1 \leqslant i, j, k, m \leqslant r).
\]

Conversely,† if (21) and (22) are satisfied, there exists a system of infinitesimal transformations satisfying relations (19), whence a transformation group with r parameters (in other words, the linear combinations of the \(X_{k}\) with constant coefficients form a Lie algebra and conversely every finite-dimensional Lie algebra can be obtained in this way).

These results are completed by studying the questions of isomorphisms. Two transformation groups are called similar if it is possible to pass from one to the other by an invertible coordinate transformation on the variables and an invertible coordinate transformation on the parameters: from the very beginning of his research, Lie had been concerned naturally with this notion when defining the ``canonical parameters''. He showed that two groups are similar if, by a transformation on the ``variables'', it is possible to transform the infinitesimal transformations of the one into those of the other ({[}4{]}, vol.~1, p.~329). A necessary condition for this to be the case, is that the Lie algebras of these two groups be isomorphic, which Lie expressed by saying that the groups are ``gleichzusammengesetzt''; but this condition is not sufficient and a whole chapter ({[}4{]}, vol.~1, Chapter 19) is devoted to obtaining supplementary conditions assuring that the groups are ``similar''. The theory of permutation groups on the other hand provided the notion of ``holohedric isomorphism'' of two such groups (isomorphism of the underlying ``abstract'' groups); Lie transposed this notion to transformation groups and showed that two such groups are ``holohedrically isomorphic'' if and only if their Lie algebras are isomorphic ({[}4{]}, vol.~1, p.~418). In particular, every transformation group is holohedrically isomorphic to each of its parameter groups and this shows that, when we wish to study the structure of the group, the ``variables'' on which it operates are of little importance and that in fact all that matters is the Lie algebra.†

Always by analogy with the theory of permutation groups, Lie introduced the notions of subgroups, normal subgroups and ``merihedric isomorphisms'' (surjective homomorphisms) and showed that they correspond to those of subalgebras, ideals and surjective Lie algebra homomorphisms; he had already previously come across an especially important example of a ``merihedric isomorphism'', the adjoint representation, and had recognized its relations with the centre of the group ({[}3f{]}, § 3, no. 9). For these results, as for the fundamental theorems, the essential tool is the Jacobi-Clebsch Theorem giving the complete integrability of a differential system (one of the forms of the theorem called ``{[}Frobenius's''); he however gave a new proof using one-parameter groups ({[}4{]}, vol.~1, Chapter 6).

The notions of transitivity and primitivity, so important for permutation groups, occurred just as naturally for ``finite continuous'' transformation groups and the Lie--Engel treatise made a detailed study of this ({[}4{]}, vol.~1, Chapter 13 and passim); the relations with the stabilizer subgroups of a point and the notion of homogeneous space were perceived (insofar as was possible without taking a global point of view) ({[}4{]}, vol.~1, p.~425).

Finally the ``dictionary'' was completed, in {[}4{]}, by the introduction of the notions of derived group and solvable group (called ``integrable group'' by Lie; this terminology, suggested by the theory of differential equations, remained in use until the works of H. Weyl) ({[}4{]}, vol.~1, p.~261 and vol.~3, pp.~678--679); the relation between commutators and brackets had elsewhere been perceived by Lie in 1883 ({[}3{]}, vol.~5, p.~358).

\subsection*{Other proofs of the fundamental theorems}

In {[}8{]} F. Schur showed that in canonical coordinates the \(\psi_{ik}\) of (15) satisfy the differential equations

\[
\frac {d}{d t} \left(t \psi_ {i k} (t a)\right) = \delta_ {i k} + \sum_ {j, l} c _ {j l} ^ {k} t a _ {l} \psi_ {i j} (t a).\tag{23}
\]

These are integrated and give a formula equivalent to the formula

\[
\varpi (\mathrm{X}) = \sum_ {n > 0} \frac {1}{(n + 1) !} (\mathrm{ad} (\mathrm{X})) ^ {n}\tag{24}
\]

of our Chapter III, §6, no. 4, Proposition 12; in particular, in canonical coordinates, the \(\psi_{ij}\) can be extended to integral functions of the \(a_k\) . F. Schur deduced a result which made precise an earlier remark of Lie: if, in definition (4) of transformation groups, it is only assumed that the \(f_i\) are of class \(\mathbf{C}^2\) , then the group is holohedrically isomorphic to an analytic group. \(\dagger\) Following his research on the integration of differential systems, E. Cartan ({[}12{]}, vol.~II \(_{2}\) , p.~371) introduced in 1904 Pfaffian forms

\[
\omega_ {k} = \sum_ {i = 1} ^ {r} \psi_ {k i} d a _ {i} (1 \leqslant i \leqslant r)\tag{25}
\]

(in the notation of (15)), called later Maurer-Cartan forms. The Maurer conditions (17) could be written

\[
d \omega_ {k} = - \frac {1}{2} \sum_ {i, j} c _ {i j} ^ {k} \omega_ {i} \wedge \omega_ {j};
\]

E. Cartan showed that the theory of finite continuous groups could be developed starting from the \(\omega_{k}\) and established the equivalence of this point of view and that of Lie. But, for him, the interest of this method was above all that it could be adapted to ``infinite continuous groups'' the theory of which he pushed much further than Lie had done and which allowed him to develop his theory of the generalized ``moving frame''.

\hnsection{THE THEORY OF LIE ALGEBRAS}

Having once acquired the correspondence between transformation groups and Lie algebras, the theory took a considerably more algebraic turn and became centred on a deep study of Lie algebras.†

A first, short period, from 1888 to 1894, marked by the works of Engel and his pupil Umlauf and especially Killing and E. Cartan, achieved a series of spectacular results on complex Lie algebras. We have seen earlier that the notion of solvable Lie algebra was due to Lie himself, who had shown (in the complex case) the theorem on the reduction of solvable linear Lie algebras to triangular form ({[}4{]}, vol.~1, p.~270). \(^{\ddagger}\) Killing observed {[}11{]} that there exists in a Lie algebra a largest solvable ideal (which we now call the radical) and that the quotient of the Lie algebra by its radical has zero radical; he called Lie algebras with zero radical semi-simple and proved that they are products of simple algebras (the latter notion having already been introduced by Lie, who had proved the simplicity of the ``classical'' algebras ({[}4{]}, vol.~3, p.~682)).

On the other hand, Killing introduced in a Lie algebra the characteristic equation \(\det(\mathrm{ad}(x)-\omega.1)=0\) , already encountered by Lie when studying 2-dimensional Lie subalgebras containing a given element of a Lie algebra. We shall return in other Historical Notes to this Book to the analysis of the methods by which Killing, whilst engaged on a penetrating study of the properties of the roots of the ``generic'' characteristic equation for a semi-simple algebra, achieved the most remarkable of his results, the complete determination of the (complex) simple Lie algebras.§

Killing proved that the derived algebra of a solvable algebra is ``of rank 0'' (which means that ad x is nilpotent for every element x of the algebra). Almost immediately, Engel showed that algebras ``of rank 0'' are solvable (this statement is essentially what we called Engel's Theorem in Chapter I, § 4, no. 2). In his thesis, E. Cartan introduced, on the other hand, what we now call the ``Killing form'' and established the two fundamental criteria which characterize by means of this form solvable Lie algebras and semi simple Lie algebras.

Killing had affirmed ({[}11{]}, IV) that the derived algebra of a Lie algebra is the sum of a semi-simple algebra and its radical, which is nilpotent, but his proof was incomplete. A little later, E. Cartan announced without proof ({[}12{]}, vol.~I₁, p.~104) that more generally every Lie algebra is the sum of its radical and a semi-simple subalgebra; the only result in this direction established indisputably at this period was a theorem of Engel affirming the existence, in every non-solvable Lie algebra, of a simple Lie subalgebra of dimension 3. The first published proof (for complex Lie algebras) of Cartan's statement is due to E. E. Levi {[}18{]}; another proof (equally valid in the real case) was given by J. H. C. Whitehead in 1936 {[}26 a{]}. In 1942, A. Malcev completed this result by the uniqueness theorem for ``Levi sections'' up to conjugation.

From his very first works, Lie was concerned with the problem of the isomorphism between any Lie algebra and a linear Lie algebra. He had believed that he had solved it affirmatively by considering the adjoint representation (and deduced from it a proof of his ``third theorem'') \([3\ d]\) ; he realized that his proof was correct only for Lie algebras with zero centre, gave another erroneous proof thereof ( \([3\ f]\) , § 3, no. 9) and then recognized ( \([3\ g]\) , p.~231) that the question remained open. In fact, it remained so for a long time and was only finally resolved affirmatively in 1935 by Ado {[}27{]}. On the other hand, Lie had essentially posed the problem of determining linear representations of simple Lie algebras of minimal dimension and had solved it for the classical groups; in his Thesis, Cartan also solved this problem for the exceptional simple algebras \(^{\dagger}\) ; the methods he used for this were to be generalized by him twenty years later to obtain all the irreducible representations of real or complex simple Lie algebras.

The property of complete reducibility of a linear representation seems to have been encountered for the first time (in a geometric form) by Study. In an unpublished manuscript, but quoted in {[}4{]}, vol.~3, pp.~785--788, he proved this property for linear representations of the Lie algebra \(\mathbf{SL}(2,\mathbf{C})\) and obtained partial results for \(\mathbf{SL}(3,\mathbf{C})\) and \(\mathbf{SL}(4,\mathbf{C})\) . Lie and Engel conjectured on this occasion that the complete reducibility theorem was true for \(\mathbf{SL}(n,\mathbf{C})\) for all \(n\) . The complete reducibility of linear representations of semi-simple Lie algebras was established by H. Weyl in 1925† by a global type of argument (see later). The first algebraic proof was obtained in 1935 by Casimir and van der Waerden {[}32{]}; other algebraic proofs have since been given by R. Brauer {[}31{]} (this is the one we have reproduced) and J. H. C. Whitehead {[}26 b{]}.

Finally, in the course of his research on the exponential mapping (cf.~infra), H. Poincaré ({[}14{]}, vol.~3) considered the associative algebra of differential operators of all orders, generated by the operators of a Lie algebra; he showed essentially that, if \((\mathbf{X}_i)_{1\leqslant i\leqslant n}\) is a basis of the Lie algebra, the associative algebra generated by the \(\mathbf{X}_i\) has as basis certain symmetric functions of the \(\mathbf{X}_i\) (sum of the non-commutative ``monomials'' derived from a given monomial by all permutations of factors). The essence of his proof was algebraic in nature and allowed him to obtain the enveloping algebra structure which we define abstractly in Chapter I. Analogous proofs were given in 1937 by G. Birkhoff {[}29 b{]} and E. Witt {[}30{]}.

Most of the results cited above are limited to real or complex Lie algebras which alone correspond to Lie groups in the usual sense. The study of Lie algebras over a field other than \(\mathbf{R}\) or \(\mathbf{C}\) was embarked upon by Jacobson {[}28 a{]} who showed that the great majority of the classical results (i.e.~those of Chapter I) remain true for any field of characteristic zero.

\hnsection{EXPONENTIAL AND HAUSDORFF FORMULA}

The first research on the exponential mapping was due to E. Study and F. Engel; Engel {[}9 b{]} remarked that the exponential is not surjective for \(\mathbf{SL}(2, \mathbf{C})\) (for example \(\begin{pmatrix} -1 & a \\ 0 & -1 \end{pmatrix}\) is not an exponential if \(a \neq 0\) ), but that it is for \(\mathbf{GL}(n, \mathbf{C})\) and hence also for \(\mathbf{PGL}(n, \mathbf{C})\) (the latter property had already been noted by Study for \(n = 2\) ); thus \(\mathbf{SL}(2, \mathbf{C})\) and \(\mathbf{PGL}(2, \mathbf{C})\) give an example of two locally isomorphic groups, which are however very different from a global point of view. Engel also showed that the exponential is surjective for the other classical groups augmented by homotheties; this work was taken up and pursued by Maurer, Study and others, without producing substantial new results.

In 1899, H. Poincaré ({[}14{]}, vol.~3, pp.~169--172 and 173--212) embarked upon the study of the exponential mapping from a different point of view. His memoirs seem to have been hastily edited, for in several places he affirmed that every element of a connected group is an exponential, whereas he gave examples to the contrary elsewhere. His results were mainly concerned with the adjoint representation: he showed that a semi-simple element of such a group G may be the exponential of an infinity of elements of the Lie algebra L(G), whereas a non-semi-simple element may not be an exponential at all. If ad(X) has no eigenvalue which is a non-zero multiple of 2πi, then exp is étale at X. He also proved that, if U and V describe loops in L(G) and W is defined by continuity such that \(e^{U} \cdot e^{V} = e^{W}\) , it is possible not to return to the original value of W. He used a residue formula which amounts essentially to

\[
\Phi (\mathrm{adX}) = \frac {1}{2 \pi i} \int \frac {\Phi (\xi) d \xi}{\xi - \mathrm{adX}}
\]

where \(\operatorname{ad}(\mathbf{X})\) is a semi-simple element whose non-zero eigenvalues are of multiplicity 1, \(\Phi\) is an integral series with sufficiently large radius of convergence and the integral is taken over a contour surrounding the eigenvalues of \(\operatorname{ad} X\) ; he also studied what happens as X tends to a transformation with multiple eigenvalues.

Research into expressions for W as a function of U and V in the formula \(e^{U}.e^{V}=e^{W}\) had already, just before Poincaré's work, been the object of two memoirs by Campbell {[}13{]}. As Baker wrote a little later ``. . . Lie theory suggests in an obvious way that the product \(e^{U}.e^{V}\) is the form \(e^{W}\) where W is a series of alternants in U and V . . .''. The later works on this subject aimed at making this assertion precise and giving an explicit formula (or a method of construction) for W (``Hausdorff formula''). After Campbell and Poincaré, Pascal, Baker {[}15{]} and Hausdorff {[}16{]} returned to the question; each considered that the proofs of his predecessors were not convincing; the principal difficulty resided in what is meant by ``alternants'': are they special elements of the Lie algebra in question, or universal ``symbolic'' expressions? Neither Campbell, nor Poincaré, nor Baker expressed himself clearly on this point. Hausdorff's memoir, on the other hand, is perfectly precise; he worked first on the algebra of (noncommutative) associative formal power series in a finite number of indeterminates and considered U, V, W as elements of this algebra. He proved the existence of W by a differential equation argument analogous to that of his predecessors. He used the same argument to prove the convergence of the series when the indeterminates are replaced by elements of a finite-dimensional Lie algebra. As Baker had remarked, and Poincaré independently, this result could be used to give a proof of Lie's third theorem; he clarified the correspondence between Lie groups and Lie algebras, for example where the commutator subgroup is concerned.

In 1947, Dynkin {[}39{]} again took up the question and obtained the explicit coefficients of the Hausdorff formula, by considering from the outset a normed Lie algebra (of finite dimension or otherwise, over R, C or an ultrametric field).†

\hnsection{LINEAR REPRESENTATIVES AND GLOBAL LIE GROUPS}

None of the works which we have just mentioned made any honest attempt to define and study global Lie groups. It was H. Weyl who made the first move in this direction. He was inspired by two theories, which until then had developed independently: that of linear representations of complex semi-simple Lie algebras, due to E. Cartan, and that of linear representations of finite groups, due to Frobenius and which had just been transposed to the orthogonal group by I. Schur, using an idea of Hurwitz. The later had shown {[}17{]} how to form invariants for the orthogonal group or unitary group by replacing the operation of the mean on a finite group by an integration with respect to an invariant measure. He had also noted that, by applying this method to the unitary group, invariants are obtained for the general linear group, the first example of the ``unitarian trick''. In 1924, I. Schur {[}20{]} used this procedure to show the complete reducibility of the representations of the orthogonal group \(\mathbf{O}(n)\) and the unitary group \(\mathbf{U}(n)\) by constructing an invariant non-degenerate positive definite Hermitian form; he deduced, by the ``unitarian trick'', the complete reducibility of the holomorphic representations of \(\mathbf{O}(n,\mathbf{C})\) and \(\mathbf{SL}(n,\mathbf{C})\) , established orthogonality relations for the characters of \(\mathbf{O}(n)\) and \(\mathbf{U}(n)\) and determined the characters of \(\mathbf{O}(n)\) . H. Weyl immediately extended this method to complex semi-simple Lie algebras {[}21{]}. Given such an algebra g, he showed that it has a ``compact real form'' (which amounts to saying that is obtained by extension of scalars from R to C from an algebra \(g_{0}\) over R whose adjoint group \(G_{0}\) is compact). Further, he showed that the fundamental group of \(G_{0}\) is finite and hence that the universal covering‡ of \(G_{0}\) is compact. He deduced, by a suitable adaptation of Schur's procedure, the complete reducibility of the representations of g and also gave, by a global method, the determination of the characters of the representations of g. In a letter to I. Schur (Sitzungsber, Berlin, 1924, pp.~338--343), H. Weyl summarized the results of Cartan not known to Schur (cf.~{[}20{]}, p.~299, note at the bottom of the page) and compared the two approaches: Cartan's method gave all the holomorphic representations of the simply connected group with Lie algebra g; in the case of the orthogonal group, representations were thus obtained of a double covering (called later the spinor group), which had escaped Schur; on the other hand, Schur's method had the advantage of showing the complete reducibility and giving the characters explicitly.

After the work of H. Weyl, E. Cartan adopted an openly global approach in his research on symmetric spaces and Lie groups. This approach was the basis of his exposition of 1930 ({[}12{]}, vol.~I₂, pp.~1165--1225) of the theory of ``finite continuous'' groups. In particular we find here the first proof of the global version of the 3rd fundamental theorem (the existence of a Lie group with given Lie algebra); Cartan also showed that every closed subgroup of a real Lie group is a Lie group (Chapter III, § 8, no. 2, Theorem 2) which generalized a result of J. von Neumann on closed subgroups of the linear group {[}23{]}. In this Memoir, von Neumann showed also that every continuous representation of a semi-simple group is real analytic.

After these works, the theory of Lie groups in the ``classical'' sense (that is in finite dimension over R or C) was more or less complete in its essentials. The first detailed exposition of it was given by Pontrjagin in his book on topological groups {[}36{]}; there he followed an approach quite close to that of Lie, but carefully distinguished between the local and the global. This was followed by Chevalley's book {[}38{]} which also contains the first systematic discussion of the theory of analytic manifolds and exterior differential calculus; the ``infinitesimal transformations'' of Lie appear there as vector fields and the Lie algebra of a Lie group G is identified with the space of left invariant fields on G. He leaves out any discussion on ``group germs'' and ``transformation groups''.

\hnsection{EXTENSIONS OF THE NOTION OF LIE GROUP}

Nowadays the vitality of Lie theory is manifested by the diversity of its applications (in topology, differential geometry, arithmetic, etc.) and by the creation of parallel theories where the underlying manifold structure is replaced by a similar structure ( \(p\) -adic manifold, algebraic variety, scheme, formal scheme, \(\ldots\) ). We are not concerned here with the history of these developments and we shall confine ourselves to those touched on in Chapter III: Banach Lie groups and \(p\) -adic Lie groups.

\textbf{(a) Banach Lie groups}

Here we are concerned with ``infinite-dimensional'' Lie groups. From the local point of view, a neighbourhood of 0 in Euclidean space is replaced by a neighbourhood of 0 in a Banach space. This is what G. Birkhoff did in 1936 {[}29 a{]}, thus achieving the notion of a complete normed Lie algebra and its correspondence with a ``group germ'' defined on an open set of a Banach space. About 1950, Dynkin completed these results by extending the Hausdorff formula to this case (cf.~supra).

The definitions and results of Birkhoff and Dynkin are local. Until recently, it does not seem that any one has sought to make the corresponding global theory explicit, no doubt because of the lack of applications.†

\textbf{(b) p-adic Lie groups}

Such groups were encountered for the first time in 1907 in the works of Hensel {[}19{]} on \(p\) -adic analytic functions (defined by expansions as integral series). He studied in particular the exponential and the logarithm; in spite of the a priori surprising behaviour of the series which define them (for example the exponential series does not converge everywhere), their fundamental functional properties remain valid, which provides a local isomorphism between the additive group and the multiplicative group of \(\mathbf{Q}_p\) (or, more generally, of any complete ultrametric field of characteristic zero).

A. Weil {[}33{]} and E. Lutz {[}34{]} were equally concerned with commutative (but this time non-linear) groups in their work on p-adic elliptic curves (1936). Besides arithmetic applications, a local isomorphism is constructed here between the group and the additive group, based on the integration of an invariant differential form. This method applies equally to Abelian varieties, as was noted by C. Chabauty soon afterwards, who used it without further explanation to prove a particular case of the ``Mordell conjecture'' {[}35{]}.

From then on, it was clear that the local theory of Lie groups could be applied with scarcely any change to the p-adic case. The fundamental theorems of the Lie groups--Lie algebras ``dictionary'' were established in 1942 in the thesis of R. Hooke {[}37{]}, a pupil of Chevalley; this work also contained the p-adic analogue of E. Cartan's theorem on closed subgroups of real Lie groups.

More recently, M. Lazard {[}42 b{]} developed a more precise form of the ``dictionary'' for compact analytic groups over \(\mathbf{Q}_{\mathrm{p}}\) . He showed that the existence of a \(p\) -adic analytic structure on a compact group G is closely related to that of certain filtrations on G and gives various applications of this (for example to the cohomology of G). One of Lazard's tools is an improvement of Dynkin's results on the convergence of the \(p\) -adic Hausdorff series {[}42 a{]}.

\hnsection{FREE LIE ALGEBRAS}

It remains to speak of a series of works on Lie algebras where the connection with the theory of Lie groups is very tenuous; this research has on the other hand important applications to the theory of ``abstract'' groups and more especially nilpotent groups.

Its origin is the work of P. Hall {[}24{]}, which appeared in 1932. However Lie algebras are not discussed here: P. Hall had in mind the study of a certain class of p-groups, those which he called ``regular''. But this led him to examine in detail iterated commutators and the lower central series of a group; on this subject he established a version of the Jacobi identity (cf.~Chapter II, § 4, no. 4, formula (20)) and the ``Hall formula''

\[
(x y) ^ {n} = x ^ {n} y ^ {n} (x, y) ^ {n (1 - n) / 2} \dots \quad (\text { cf.   Chapter   II }, \S 5, \text { Exercise   9 }).
\]

Almost immediately (in 1935--1937) appeared the fundamental works of W. Magnus ({[}25 a{]} and {[}25 b{]}) and E. Witt {[}30{]}. In {[}25 a{]} Magnus used the same algebra of formal power series \(\hat{\mathbf{A}}\) as Hausdorff (since called ``Magnus algebra''); he embedded the free group F in it and used the natural filtration of \(\hat{\mathbf{A}}\) to obtain a decreasing sequence \((\mathbf{F}_{n})\) of subgroups of F; it is one of the first examples of a filtration. He conjectured that the \(F_{n}\) coincide with the terms of the lower central series of F. This conjecture was proved in his second memoir {[}25 b{]}; also in this work he showed explicitly the close relation between his ideas and those of P. Hall and defined the free Lie algebra L (as a subalgebra of \(\hat{\mathbf{A}}\)) which he showed essentially to be identified with the graded algebra of F. In {[}30{]}, Witt completed this result on various points. He showed notably that the enveloping algebra of L is a free associative algebra and deduced immediately the rank of the homogeneous components of L (``Witt formulae'').

As for the proof of the basis of L known as the ``Hall basis'' (cf.~Chapter II, § 2, no. 11), it seems that it only appeared in 1950 in a note of M. Hall {[}40{]}, although it was implicit in the works of P. Hall and W. Magnus quoted above.

\specialchapter{Bibliography}

\begin{enumerate}
\def\labelenumi{\arabic{enumi}.}
\item
  F. KLEIN and S. LIE: (a) Sur une certaine famille de courbes et surfaces, C. R. Acad. Sci., 70 (1870), pp.~1222--1226 and 1275--1279 (={[}2{]}, pp.~416--420 and {[}3{]}, vol.~1, pp.~78--85); (b) Über diejenigen ebenen Kurven, welche durch ein geschlossenes System von einfach unendlich vielen vertauschbaren linearen Transformationen in sich übergehen, Math. Ann., 4 (1871), pp.~50--84 (={[}2{]}, pp.~424--459 and {[}3{]}, vol.~1, Abh. XIV, pp.~229--266).
\item
  F. KLEIN, Gesammelte mathematische Abhandlungen, Bd. I, Berlin (Springer), 1921.
\item
  S. LIE, Gesammelte Abhandlungen, 7 vol., Leipzig (Teubner): (a) Über die Reziprozitätsverhältnisse des Reyeschen Komplexes, vol.~I, Abh. V, pp.~68--77 (=Gött. Nach. (1870), pp.~53--66); (b) Über eine Klasse geometrischer Transformationen, vol.~I, Abh. XII, pp.~153--214 (=Christiana For. (1871), pp.~182--245); (c) Über partielle Differentialgleichungen erster Ordnung, vol.~III, Abh. VII, pp.~32--63 (=Christiana For. (1873), pp.~16--51); (d) Theorie der Transformationsgruppen II, vol.~V, Abh. III, pp.~42--75 (=Archiv. Math., I (1876), pp.~152--193); (e) Ueber Gruppen von Transformationen, vol.~V, Abh. I, pp.~1--8 (=Gött. Nachr. (1874), pp.~529--542); (f) Allgemeine Untersuchungen über Differentialgleichungen, die eine kontinuierliche endliche Gruppe gestatten, vol.~VI, Abh. III, pp.~139--223 (=Math. Ann., 25 (1885), pp.~71--151); (g) Beiträge zur allgemeinen Transformationenstheorie, vol.~VI, Abh. V, pp.~230--236 (=Leipziger Ber. (1888), pp.~14--21); (h) Theorie der Transformationsgruppen III, vol.~V, Abh. IV, pp.~78--133 (=Archiv. Math., vol.~III, (1887), pp.~93--165).
\item
  S. LIE and F. ENGEL, Theorie der Transformationsgruppen, 3 vol., Leipzig (Teubner), 1888--1893.
\item
  J. J. SYLVESTER, Collected Mathematical Papers, 4 vols., Cambridge, 1904-1911.
\item
  A. CAYLEY, Collected Mathematical Papers, 13 vols., Cambridge, 1889-1898.
\item
  C. JORDAN, Mémoire sur les groupes de mouvements, Annali di Math., 11 (1868--1869), pp.~167--215 and 332--345 (= Oeuvres, vol.~IV, pp.~231--302).
\item
  F. SCHUR: (a) Zur Theorie der aus Haupteinheiten gebildeten Komplexen, Math. Ann., 33 (1889), pp.~49--60; (b) Neue Begründung der Theorie der endlichen Transformationsgruppen, Math. Ann., 35 (1890), pp.~161--197; (c) Zur Theorie der endlichen Transformationsgruppen, Math. Ann., 38 (1891), pp.~273--286; (d) Über den analytischen Character der eine endliche continuierliche Transformationsgruppe darstellende Funktionen, Math. Ann., 41 (1893), pp.~509--538.
\item
  F. ENGEL: (a) Über die Definitionsgleichung der continuerlichen Transformationsgruppen, Math. Ann., 27 (1886), pp.~1--57; (b) Die Erzeugung der endlichen Transformationen einer projektiven Gruppe durch die infinitesimalen Transformationen der Gruppe, I, Leipzig Ber., 44 (1892), pp.~279--296, II (mit Beiträgen von E. Study), ibid., 45 (1893), pp.~659--696.
\item
  L. MAURER, Über allgemeinere Invarianten-Systeme, Sitzungsber. München, 18 (1888), pp.~103--150.
\item
  W. KILLING, Die Zusammensetzung der stetigen endlichen Transformationsgruppen: (I) Math. Ann., 31 (1888), pp.~252--290; (II) ibid., 33 (1889), pp.~1--48; (III) ibid., 34 (1889), pp.~57--122; (IV) ibid., 36 (1890), pp.~161--189.
\item
  E. CARTAN, Oeuvres complètes, 6 vol., Paris (Gauthier-Villars), 1952--54.
\item
  J. E. CAMPBELL: (a) On a law of combination of operators bearing on the theory of continuous transformation groups, Proc. London Math. Soc., (1) 28 (1897), pp.~381--390; (b) On a law of combination of operators (second paper), ibid., 29 (1898), pp.~14--32.
\item
  H. POINCARÉ, Oeuvres, 11 vol., Paris (Gauthier-Villars), 1916--1956.
\item
  H. F. BAKER, Alternants and continuous groups, Proc. London Math. Soc., (2) 3 (1905), pp.~24--47.
\item
  F. HAUSDORFF, Die symbolische Exponentialformel in der Gruppentheorie, Leipziger Ber., 58 (1906), pp.~19--48.
\item
  A. HURWITZ, Über die Erzeugung der Invarianten durch Integration, Gött, Nachr. (1897), pp.~71--90 (=Math. Werke, vol.~II, pp.~546--564).
\item
  E. E. LEVI, Sulla struttura dei gruppi finiti e continui, Atti Acc. Sci. Torino, 40 (1905), pp.~551--565 (= Opere, vol.~I, pp.~101--115).
\item
  K. HENSEL, Über die arithmetischen Eigenschaften der Zahlen, Jahresber. der D.M.V., 16 (1907), pp.~299--319, 388--393, 474--496.
\item
  I. SCHUR, Neue Anwendungen der Integralrechnung auf Probleme der Invariantentheorie, Sitzungsber. Berlin, 1924, pp.~189--208, 297--321, 346--355.
\item
  H. WEYL, Theorie der Darstellung kontinuierlicher halb-einfacher Gruppen durch lineare Transformationen, I, Math. Zeitschr., 23, (1925),
\end{enumerate}

pp.~271--309; II, ibid., 24 (1926), pp.~328--376; III, ibid., 24 (1926), pp.~377--395 (=Werke, vol.~II, pp.~543--647).

\begin{enumerate}
\def\labelenumi{\arabic{enumi}.}
\setcounter{enumi}{21}
\item
  O. SCHREIER: (a) Abstrakte kontinuierliche Gruppen, Abh. math. Sem. Hamburg, 4 (1926), pp.~15--32; (b) Die Verwandschaft stetiger Gruppen in grossen, ibid., 5 (1927), pp.~233--244.
\item
  J. VON NEUMANN, Zur Theorie der Darstellung kontinuierlicher Gruppen, Sitzungsber. Berlin, 1927, pp.~76--90 (=Collected Works, vol.~I, pp.~134--148).
\item
  P. HALL, A contribution to the theory of groups of prime power order, Proc. London Math. Soc., (3) 4 (1932), pp.~29--95.
\item
  W. MAGNUS: (a) Beziehungen zwischen Gruppen und Idealen in einem speziellen Ring, Math. Ann., 111 (1935), pp.~259--280; (b) Über Beziehungen zwischen höheren Kommutatoren, J. Crelle, 177 (1937), pp.~105--115.
\item
  J. H. C. WHITEHEAD: (a) On the decomposition of an infinitesimal group, Proc. Camb. Phil. Soc., 32 (1936), pp.~229--237 (=Mathematical Works, I, pp.~281--289); (b) Certain equations in the algebra of a semi-simple infinitesimal group, Quart. Journ. of Math., (2) 8 (1937), pp.~220--237 (=Mathematical Works, I, pp.~291--308).
\item
  I. Ado: (a) Note on the representation of finite continuous groups by means of linear substitutions (in Russian), Bull. Phys. Math. Soc. Kazan, 7 (1935), pp.~3--43; (b) The representation of Lie algebras by matrices (in Russian), Uspehi Mat. Nauk, 2 (1947), pp.~159--173 (English translation: Amer. Math. Soc. Transl., (1) 9, pp.~308--327).
\item
  N. JACOBSON: (a) Rational methods in the theory of Lie algebras, Ann. of Math., 36 (1935), pp.~875--881; (b) Classes of restricted Lie algebras of characteristic p, II, Duke Math. Journal, 10 (1943), pp.~107--121.
\item
  G. BIRKHOFF: (a) Continuous groups and linear spaces; Rec. Math. Moscou, 1 (1936), pp.~635--642; (b) Representability of Lie algebras and Lie groups by matrices, Ann. of Math., 38 (1937), pp.~526--532.
\item
  E. WITT, Treue Darstellung Lieschen Ringe, J. Creille, 177 (1937), pp.~152--160.
\item
  R. BRAUER, Eine Bedingung für vollständige Reduzibilität von Darstellungen gewöhnlicher und infinitesimaler Gruppen, Math. Zeitschr., 41 (1936), pp.~330--339.
\item
  H. CASIMIR-B. L. VAN DER WAERDEN, Algebraischer Beweis der vollständigen Reduzibilität der Darstellungen halbeinfacher Liescher Gruppen, Math. Ann., 111 (1935), pp.~1--12.
\item
  A. WEIL, Sur les fonctions elliptiques p-adiques, C. R. Acad. Sci., 203 (1935), p.~22.
\item
  E. Lutz, Sur l'équation \(y^{2} = x^{3} - Ax - B\) dans les corps \(p\) -adiques, \(J\) . Crelle, 177 (1937), pp.~237-247.
\item
  C. CHABAUTY, Sur les points rationels des courbes algébriques de genre supérieur à l'unité, C. R. Acad. Sci., 212 (1941), pp.~882--884.
\end{enumerate}

\begin{enumerate}
\def\labelenumi{\arabic{enumi}.}
\setcounter{enumi}{35}
\item
  L. S. PONTRJAGIN, Topological groups, Princeton Univ. Press, 1939.
\item
  R. HOOKE, Linear \(p\) -adic groups and their Lie algebras, Ann. of Math., 43 (1942), pp.~641-655.
\item
  C. CHEVALLEY, Theory of Lie groups, Princeton Univ. Press, 1946.
\item
  E. DYNKIN: (a) Evaluation of the coefficients of the Campbell-Hausdorff formula (in Russian), Dokl. Akad. Nauk, 57 (1947), pp.~323-326; (b) Normed Lie algebras and analytic groups (in Russian), Uspehi Mat. Nauk, 5 (1950), pp.~135-186 (English translation: Amer. Math. Soc. Transl., (1) 9, pp.~470-534).
\item
  M. HALL, A basis for free Lie rings and higher commutators in free groups, Proc. Amer. Math. Soc., 1 (1950), pp.~575--581.
\item
  D. MONTGOMERY-L. ZIPPIN, Topological Transformation Groups, New York (Interscience), 1955.
\item
  M. LAZARD: (a) Quelques calculs concernant la formule de Haudorff, Bull. Soc. Math. France, 91 (1963), pp.~435--451; (b) Groupes analytiques p-adiques, Publ. Math. I.H.E.S., no. 26 (1965), pp.~389--603.
\end{enumerate}

\(\mathcal{C}_k\mathfrak{g}\) (g a Lie algebra): I.1.6.

ad \(_{g}\) x, ad x (x an element of a Lie algebra g): I.1.2.

\(g_{(K_{1})}\) (g a Lie algebra): I.1.9.

\(e^u\) , \(\exp u\) ( \(u\) a nilpotent endomorphism of a vector space over a field of characteristic 0): I.6.8.

\(\mathcal{D}\mathfrak{g}, \mathcal{D}^{\kappa}\mathfrak{g}, \mathcal{C}^{\kappa}\mathfrak{g}\) (g a Lie algebra): I.1.5.
